% -----------------------------------------------------------
\section{Talks}
% -----------------------------------------------------------
	\label{TALKS}
	
	% ----------------------
	\subsection{15 February 2026: Sunday school on Moses 5--8}
	% ----------------------
		\label{TALKS-15-02-2026}
		
		\begin{compactitem}
		
			\item Hi everyone.
				\begin{compactitem}
					\item sdf
				\end{compactitem}
				
			\item Today we're going to be thinking about the broad sweep of some of the events happening in the book of Moses.
				\begin{compactitem}
					\item We'll be focusing especially on issues relating to the priesthood.
					\item I want to try to bring out the importance of the priesthood in our connection to Christ, and especially to the atonement.
					\item That's the plan, at least. We'll see what we're able to accomplish.
				\end{compactitem}
				
			\item Let's begin by thinking about the priesthood in general.
				\begin{compactitem}
					\item Open to \hyperref[DC107-1]{DC 107:1}.
					\item This section is about many things, but in the 1835 and 1844 DC versions of this revelation, there is a huge heading at the beginning which reads, ``ON PRIESTHOOD.''
					\item So it's about the priesthood.
					\item We're going to read the first four verses.
				\end{compactitem}
				
			% Actual order
				% Read DC 107:1--4 about the name of the priesthood
				% Read DC 76:57 about the ``order of Enoch''
				% ``I want to think more about the idea of an `order.' Two mistakes (?): we associate it with males; we associate it with specific offices or duties. But both of those things are inessential to the actual operation of the priesthood, and we, all of us, enter the order''
				% Read JST Genesis 14 to get more about this
				% 
				
			% PRETTY SURE THIS IS THE ENTRANCE POINT INTO THE LESSON
			\item The ``order of Enoch''
				\begin{compactitem}
					\item An ``order'' having to do with societal transformation and physical translation (?)
					\item The result is not just individual spiritual transformation but group transformation
					\item \hyperref[DC76-57]{DC 76:57}, the ``order of Enoch''
					\item \hyperref[JST-OTR1-GENESIS-14]{JST Genesis 14}
				\end{compactitem}
				
			\item Unity versus degrees of the priesthood
				\begin{compactitem}
					\item JS quotes about ``all priesthood is Melchizedek''
					\item \hyperlink{EMS-1841-JS-5-JAN-1841-WC-b}{5 January 1841 remarks} on this
					\item ``Order'' may refer to \textit{our} capacity to do things.
					\item Analogy to Newtonian and Einsteinian physics, quantum mechanics
					\item \hyperref[DC107-1]{DC 107:1--4}
				\end{compactitem}
				
			\item Entering the ``order''
				\begin{compactitem}
					\item It is not an administrative status, and it has nothing to do with an ``office'' or some other kind of ``stewardship'' in the priesthood.
					\item It is a spiritual status describing your membership in the church of the Firstborn.
					\item You enter the order by the spiritual rebirth of baptism and the Holy Ghost.
						\begin{compactitem}
							\item \hyperref[MOSES-6-64]{Moses 6:64--68}
							\item \hyperref[HEBREWS-12-23]{Hebrews 12:23}
							\item \hyperref[MOSIAH-5-7]{Mosiah 5:7}
						\end{compactitem}
				\end{compactitem}
				
			\item Orders of the priesthood as orders of personal transformation through the power of Christ
				\begin{compactitem}
					\item Enoch's transformation in the latter half of Moses 7?
					\item Here we make the connection between the priesthood and the atonement?
						\begin{compactitem}
							\item Each priesthood ordinance is a ``similitude'' that connects us to the atonement.
							\item \hyperref[MOSES-5-5]{Moses 5:5--9}
							\item The atonement is the power source; the priesthood, or the system of laws governing physical interactions, is a system intended to continually increase our access to that power.
							\item \hyperref[MOSES-6-59]{Moses 6:59--60}
							\item \hyperref[ALMA-13]{Alma 13} (?)
							\item BY quotes on the priesthood?
						\end{compactitem}
				\end{compactitem}
		
		\end{compactitem}
	
	% ----------------------
	\subsection{18 January 2026: Sunday school on Moses 1, Abraham 3}
	% ----------------------
		\label{TALKS-18-01-2026}
		
		\begin{compactitem}
		
			\item Hi everyone.
				\begin{compactitem}
					\item sdf
				\end{compactitem} 
				
			\item First of all, let's begin by talking about the Book of Moses.
				\begin{compactitem}
					\item What is this book, exactly?
					\item Moses chapter 1 exists in at least three distinct manuscript versions from around the time it was originally produced.
						\begin{compactitem}
							\item These three are Old Testament Revision 1, the John Whitmer first copy of the Old Testament revision, and Old Testament Revision 2.
							\item In those versions it is described as ``A Revelation given to Joseph the Revelator June 1830.''
								\begin{compactitem}
									\item OTR1: ``A Revelation given to Joseph the Revelator June 1830''
									\item OTRJW: ``A Revelation given to Joseph the Revelator June 1830''
									\item OTR2: ``A Revelation, given to Joseph the Seer, June, 1830.''
									\item (It is also described this way in various journal copies, but presumably they're just copying it from the manuscripts.)
								\end{compactitem}
						\end{compactitem}
					\item \tr{Why would it be significant that this text in particular, what we now call ``Moses 1,'' is described as a \textit{revelation}?}
						\begin{compactitem}
							\item Does this change how we read it, or how we think about what it says?
							\item How is it being a ``revelation'' different from being a ``translation,'' or something like that?
						\end{compactitem}
				\end{compactitem}
				
			\item Roughly speaking, we can divide the chapter into three different sections.
				\begin{compactitem}
					\item The first section covers verses 1--11.
						\begin{compactitem}
							\item This is the section where Moses has some initial experiences with God.
						\end{compactitem}
					\item The second section covers verses 12--23.
						\begin{compactitem}
							\item This is the section where Moses has his encounter with Satan.
						\end{compactitem}
				\end{compactitem}
				
			\item Section 01, Moses 1:1--11
				\begin{compactitem}
					\item 
					\item Repetition of ``son''-relation
						\begin{compactitem}
							\item 11 mentions (that I counted, I could be wrong) of a ``son''- or child-relation to God.
							\item 4 in the first section.
							\item 6 in the second section.
							\item 1 in the third section.
							\item First is verse 4.
							\item Then we get two in verse 6.
							\item \tr{What does it mean that Moses is ``in the similitude of mine Only Begotten'' (verse 6)? Is this something additional to being God's ``son,'' or is it just another way of saying the same thing?}
						\end{compactitem}
					\item In verse 9, ``the presence of God withdrew from Moses, that his glory was not upon'' him.
						\begin{compactitem}
							\item \tr{Has anyone ever experienced something like this before, like a tangible withdrawal of spiritual presence in your life?}
							\item Notice that this withdrawal of ``presence'' and ``glory'' has nothing to do with \textit{sin}, or whether Moses was righteous or not.
							\item In the course of our experiences, there are just natural movements of the Holy Ghost in us, totally unconnected to whether we are living righteously or not. 
							\item It's something like the tides of the ocean: it's sometimes in, it's sometimes out, as the result of natural processes at work in us.
							\item (\hyperref[DC76-116]{DC 76:116--118} gives us a description of a general principle involved here.)
							\item (Something similar happens to Lamoni in the Book of Mormon, and Alma, and so on.)
						\end{compactitem}
					\item Verse 10, the conclusion Moses reaches after this experience.
					\item Verse 11, the difference between ``natural'' and ``spiritual'' eyes.
						\begin{compactitem}
							\item \tr{Have you ever had experiences where you felt like your perception was quickened on something like this way, and you could see things you otherwise wouldn't be able to?}
							\item Seeing sometihng with ``spiritual eyes'' is not hallucinating; it's just a changed form of perception that is possible through the action of the Spirit on us.
							\item We'll see more of this later in the chapter.
						\end{compactitem}
					\item Notice that these are all \textit{physical} effects that are happening to Moses's \textit{body}.
						\begin{compactitem}
							\item God's ``glory,'' the ``Spirit,'' God's ``love,'' his ``power,'' and so on, all these terms refer to specific physical phenomena that really exist, and really have actual effects on our bodies and on our spirits.
							\item They are \textit{all} real.
							\item We don't know the physical mechanisms yet, and as we learn more, we will undoubtedly be wrong about in what we think about those mechanisms.
							\item The same thing happened in the history of science with gravity, electricity, and other phenomena that had detectable physical effects, but seemed mysterious.
							\item As our knowledge increases, all of these events will fall into place as comprehensible, physical phenomena that produce these tangible effects in us.
						\end{compactitem}
					\item \textbf{Moses is invited to learn the fact that there is more to reality than we are able to perceive or experience in normal circumstances. Being a child of God, or that relationship, is part of that additional reality.}
						\begin{compactitem}
							\item And there must be some physical change over our bodies that happens in order to make this additional perception possible (in some clear way---we perceive it dimly in other ways).
						\end{compactitem}
				\end{compactitem}
				
			\item Section 02, Moses 1:12--23
				\begin{compactitem}
					\item Now Satan shows up.
						\begin{compactitem}
							\item Although it feels like an extended conversation, it actually isn't. 
							\item According to the text, Satan actually only speaks twice, and he says something similar each time.
							\item In verse 12 he says, ``Son of man, worship me.''
							\item And then in verse 19 he says, ``I am the Only Begotten, worship me.''
							\item So his message to Moses, and his message to all of you today and throughout all the rest of time, is two things: first, you \textit{don't} actually have that relationship to God that you thought you did; and second, therefore, worship me (and this ``worship'' can take any form you want, by the way.)
							\item \tr{Why would Satan say that he is the ``Only Begotten''?}
								\begin{compactitem}
									\item \hyperref[MOSES-4]{Moses 4} has some helpful material here.
									\item \hyperref[MOSES-6-65]{Moses 6:65--68} are crucial here too.
									\item Satan now stands outside the God - Christ - Us lineal relationship, and there's apparently no way back in.
								\end{compactitem}
							\item In saying ``I am the Only Begotten,'' Satan may also be referring to God's statement from before, where Moses is in the \textit{similitude} of the Only Begotten---Satan is trying to say, no, you're not in \textit{Christ's} lineage, you're in \textit{mine}.
						\end{compactitem}
					\item Satan's rant (verses 19--20)
						\begin{compactitem}
							\item Moses begins to ``fear exceedingly.''
							\item Then he ``saw the bitterness of hell.''
							\item Is it that the fear \textit{leads} to bitterness? Or what is the relationship between Moses's fearing, and the bitterness? What \textit{is} the bitterness of hell, anyway?
							\item Let me suggest that, just as Moses's perception was changed by the presence of God, such that he could see things that he was not able to see before, in the very same way, \textit{fear} changes our perception too.
							\item Fear reveals to us things that we would have never seen before---it reveals to us possibilities that we couldn't otherwise perceive.
							\item The fear here is not a bad thing, I think. It's a tool that shows Moses something---the other choice, the second of the fundamental choice that everyone on earth, in one way or another, must make.
							\item And that choice is between Christ and Satan---or between light, goodness, growth, beauty, and power, or the decay, diminishment, and darkness, and dissolution.
						\end{compactitem}
					\item Moses learns this lesson for himself, and makes the choice. 
						\begin{compactitem}
							\item It is not until Moses commands Satan in the name of Christ that Satan actually departs.
							\item Moses has learned for himself the source of power that will dispel those things.
							\item Satan departs, and Moses could not see him anymore, even with his natural eyes.
						\end{compactitem} 
					\item \textbf{In this section, Moses learns that there are many powers in the natural world, but because of his experiences with God, he already knows that there is more to reality than that natural world.}
				\end{compactitem}
				
			\item Section 03, Moses 1:24--42
				\begin{compactitem}
					\item God's glory comes again. 
					\item Notice the comparison between verses 8 and 27--28.
						\begin{compactitem}
							\item In verse 8 he sees the ``world'' and ``the children of men.''
							\item But in verse 27 he sees the ``earth,'' but his perception of it is different this time. This time he sees ``all of it,'' and apparently we're supposed to understand that very literally, because he sees every single ``particle'' of it. 
							item And then in verse 28 he sees all the people living on it, but he sees their \textit{souls}, or sees them \textit{as} souls.
							\item So he sees the same things as before, at least in this portion of his experience, \textit{but his perception has been changed, and he is able to see and understand them in new ways}.
							\item This is characteristic of righteousness. Just as God's glory can have a physical effect on us, righteousness also has a physical effect on us. 
							\item Moses has now been changed.
							\item (\hyperref[DC76-94]{DC 76:94}, ``see as we are seen'')
						\end{compactitem}
					\item Moses learns for himself that yes, all of this is \textit{real}.
					\item And he then is invited to know for himself much, much more.
						\begin{compactitem}
							\item In verse 18 Moses says, ``I have other things to inquire of him.''
							\item And he gets that chance later.
							\item We will get that chance too; that is at the heart of our discipleship of Christ.
						\end{compactitem} 
				\end{compactitem}
			
			\item Moses's experience as a representation of the plan of salvation
				\begin{compactitem}
					\item In God's presence, something like a fall, loss of physical strength, encounter with Satan, chooses not to worship Satan, re-enters God's presence with greater views than what he had before.
					\item This then gets expounded on in the next six chapters of Moses, or Moses 2--7.
					\item I suggest that we might more profitably read all of these as re-telling the story of the plan of salvation in a different way, with a different perspective, inspired by some passages in the Bible.
						\begin{compactitem}
							\item But it's all a revelation in the end.
						\end{compactitem}
				\end{compactitem}
				
			\item Scriptures I might want
				\begin{compactitem}
					\item \hyperref[MOSES-4-1]{Moses 4:1--4}, where Satan wants to be God's ``son''
					\item \hyperref[ALMA-32-35]{Alma 32:35}, about how what we come to know is ``light''
					\item \hyperref[MORONI-7]{Moroni 7}, on making judgments
				\end{compactitem}
			
			\item Supplementary materials
				\begin{compactitem}
					\item The JST translation timeline in the JST book
					\item Alternate versions of Moses 1:39
					\item Timeline of the Book of Abraham (if I can find one)
					\item Comment on JS's interpretation of Genesis 1:1 in King Follet
					\item Comment on JS's ``without form and void'' thing, with Hebrew grammar
				\end{compactitem}
		
		\end{compactitem}
	
	% ----------------------
	\subsection{7 December 2025: Sunday school on DC 137--138}
	% ----------------------
		\label{TALKS-07-12-2025}
		
		\begin{compactitem}

			\item Hi everyone.
				\begin{compactitem}
					\item Today we're talking about DC 137 and 138.
					\item Both of these sections are visions had and reported by prophets.
					\item The first actually takes us back in time to the Kirtland Temple in 1836.
					\item The second then takes us very \textit{forward} in time to 1918.
					\item So these visions are separated from each other by about a hundred years.
					\item And I'm actually going to throw a third vision into the mix as well, one had by David O. McKay when he was a young missionary in the Pacific Islands. Some of you might already be familiar with it.
					\item \tr{Make sure to pass out copies of McKay's vision.}
				\end{compactitem}
				
			\item So we're going to talk about three visions today. 
				\begin{compactitem}
					\item Each was given in very different circumstances, but what they have in common is a focus on the celestial world---on Christ, as the one who makes our return to the celestial world possible, and on the people who end up entering into that celestial world.
					\item We're first going to talk about the circumstances that brought these visions about, or the context in which they were received.
					\item We're going to see if we can find any patterns among them, in order to understand, if possible, some takeaways for the ways that we receive revelation.
					\item Then we're going to look at what the three visions say about people---about us, as part of the large group of people who want to return to the presence of our heavenly parents.
					\item We'll go in chronological order, so that means we'll start with DC 137.
				\end{compactitem}
				
			\item Context for DC 137
				\begin{compactitem}
					\item Let's get started by talking about the context for DC 137. 
					\item It's a vision had by Joseph Smith in 1836. \tr{But what is the context? What is happening, and what prompts this vision?}
					\item JS and a small group of other people are in the Kirtland Temple, beginning to perform the ordinances that constituted the ``endowment'' at that time. 
					\item (See \hyperlink{EMS-1836-JS-21-JAN-1836-h}{here} for the JS journal entry.)
					\item JS has just finished blessing his own father (the presidency also blessed him).
					\item Then his father blesses him in turn (and the presidency also blesses JS).
					\item The journal entry says, ``all of the presidency laid their hands upon me and pronounced upon my head many prophesies, and blessings, many of which I shall not notice at this time, but as Paul said, so say I, let us come to vissions and revelations.''
					\item Then immediately after that you get a ``\sout{the},'' struck out, then ``The heavens were opened upon us and I beheld...''
					\item There is also a lot of material that comes \textit{after} the end of section 137.
					\item It goes on to talk about the twelve apostles, William McLellin, Brigham Young, and then the twelve apostles again; it ends this way: ``I also beheld the redemption of Zion, and many things which the toung of man, cannot discribe in full.''
					\item And beyond that, the next day, on January 22, there were more anointings and similar spiritual manifestations.
					\item Then it happens again on January 28.
					\item Then again on February 6.
					\item And there could be other dates in between all these days that we just don't have clear records for.
					\item In the entry at the end of January 28, JS's journal says, ``after these quorems were dismissed I retired to my home filled with the spirit \& my soul cried hossannah to God \& the Lamb through <the> silent watches of the night \& while my eyes were closed in sleep the visions of the Lord were sweet unto me \& his glory was round about me praise the Lord.''
					\item \tr{So what is the specific context that brings about this vision?}
						\begin{compactitem}
							\item \textbf{performing rituals for others $\rightarrow$ receives rituals from others $\rightarrow$ receives blessings from others $\rightarrow$ vision}
						\end{compactitem}
				\end{compactitem}
				
			\item Context for DC 138
				\begin{compactitem}
					\item Now let's turn to section 138. \tr{What is the context for this vision? What happens to bring it about?}
					\item He's in his room; seems to evening, but the text doesn't actually say that.
					\item He's thinking about the scriptures; he's thinking about the atonement.
					\item \tr{So what is the specific context that brings about this vision?}
						\begin{compactitem}
							\item pondering on the scriptures $\rightarrow$ reflecting on atonement and love of God $\rightarrow$ reads the scriptures $\rightarrow$ pondering $\rightarrow$ eyes of understanding opened/spirit arrives $\rightarrow$ vision
						\end{compactitem} 
				\end{compactitem}
				
			\item McKay's vision, and the context for it
				\begin{compactitem}
					\item He's sailing, and seeing various sights near land.
					\item \tr{The context and chain of events:}
						\begin{compactitem}
							\item beautiful sights in nature/stunning sunset and reflection on the water $\rightarrow$ pondering the beauty of what he's seen $\rightarrow$ compares to the beautiful and sublime character of friends $\rightarrow$ falls asleep $\rightarrow$ vision
						\end{compactitem} 
					\item Link to the vision: \href{https://drive.google.com/file/d/1IedEX3011opw9AMbH9sm2vDDDlwouLoh/view?usp=sharing}{here}.
				\end{compactitem}
				
			\item Patterns in these three visions
				\begin{compactitem}
					\item All three involve an element of what we might call ``pondering.''
						\begin{compactitem}
							\item That word isn't used in all three accounts, but all involve the receiver of the vision reflecting on what they are experiencing, and letting it settle on them. They don't interrogate it, they just experience it, and then let themselves be guided by what the experience brings up inside themselves.
							\item They also do a lot of pondering \textit{before} the experience, and not only immediately before, but also \textit{regularly}. The pondering on spiritual matters is just part of what they do, or part of who they are.
							\item This is part of what the scriptures mean by being ``spiritually minded.''
							\item Let's read from \hyperref[ROMANS-8-5]{Romans 8:5--6}.
							\item In verse 5, the language used is a way of talking in ancient Greek that refers to ``favoring or siding with a particular political party or cause.'' Paul seems to have meant to ``communicate that the pattern of life dictated by the flesh and the pattern of life dictated by the Spirit are rivals grappling for influence over the person'' (Romans ZECNT PDF 691).
							\item It is this entire \textit{pattern} of life that is then referred to in verse 6. This pattern involves our thought, desire, and action, as well as anything else that could be part of our mental nature.
							\item Other translations I have seen for this verse include:
								\begin{compactitem}
									\item ZECNT: ``For the mindset of the flesh is death, but the mindset of the Spirit is life and peace.''
									\item NIGTC: ``The mind controlled by the sinful nature is death, but the mind controlled by the Spirit is life and peace.''
									\item WBC: ``For the flesh's way of thinking is death, whereas the Spirit's way of thinking is life and peace.''
								\end{compactitem}
							\item We need to have this kind of \textit{mindset}, or \textit{way of thinking}, or mind \textit{controlled} by spiritual matters, in order to reap the benefits---some of which are, sometimes, penetrating visions.
						\end{compactitem}
					\item All three are experiencing powerful feelings that come from special human experiences.
						\begin{compactitem}
							\item JS is feeling love and fellowship and connection with others.
							\item F. Smith is thinking about his connection to God, and God's connection to us.
							\item McKay is thinking about nature.
							\item These three connections are very important: to other people, to God, and to the natural world.
							\item In all three cases, these connections were unmediated, meaning, they were experienced directly. And they had a powerful impact because these connections are the deepest kinds of connections that human beings can experience: with each other, with God, and with the world.
						\end{compactitem}
					\item All three involve a ``passive'' element, two of them grammatically so:
						\begin{compactitem}
							\item JS's vision/137: ``The heavens were opened upon us''
							\item 138: ``the eyes of my understanding were opened''
							\item McKay: no passive voice, but he falls asleep and then sees the things in the vision
							\item In other words, they were \textit{acted on} in receiving this vision; 
							\item It is also significant that \textit{none of them asked to receive this vision}. \tr{Why would this be significant?}
						\end{compactitem}
					\item \tr{What do these patterns teach us about revelation in our own lives?}
				\end{compactitem}
		
			\item Let's now get into the visions a little bit, to see what kind of people are the ones that were seen by the receivers of these visions.
				\begin{compactitem}
					\item First, let's go to DC 138. 
					\item In verse 11 he reports that he saw ``the hosts of the dead, both small and great.'' This seems to be something he sees from afar, or it's a description of something that he takes in from a broad perspective.
					\item Because then in the next verse, he says that there were also ``gathered together in one place'' an enormous number of the ``spirits of the just.'' So in seeing these spirits he seems to be paying attention to a single part or place in what he had seen before.
					\item This description of the ``spirits of the just'' comes from \hyperref[HEBREWS-12-23]{Hebrews 12:23}, where we read about people ascending to God and ``and to the spirits of just men made perfect.''
					\item The language then gets used again in \hyperref[DC76-69]{DC 76:69}, in a description of the people in the celestial kingdom. There we read, ``These are they who are just men made perfect through Jesus the mediator of the new covenant, who wrought out this perfect atonement through the shedding of his own blood.''
					\item So in that description, it's made clear that these people who are ``just'' become something \textit{more} than ``just'' through Jesus, and through his atonement. They become ``perfect.''
					\item But now if we look again at section 138, we get some additional insight about the people who were gathered together in this place, waiting for Christ to appear after his death.
					\item First we read, continuing in verse 12, that ``they had been faithful in the testimony of Jesus while they lived in mortality.''
					\item Now go into the first part of 13. These were people who ``had offered sacrifice in the \hyperref[SIMILITUDE]{similitude} of the great sacrifice of the Son of God.''
					\item \textit{They themselves had offered sacrifice}; in 76 we read that it was the atonement that took these spirits of the just, and made them perfect. 
					\item But part of their journey involved they themselves offering sacrifice, and their particular sacrifice was ``in the similitude of the great sacrifice'' that Christ made.
					\item Let's talk about this a bit.
					\item \tr{First, what does it mean to say that their sacrifice was ``in the similitude'' of Christ's great sacrifice?}
					\item There are some verses of scripture that may help us break down the idea of a sacrifice being made ``in the similitude'' of Christ's great sacrifice.
						\begin{compactitem}
							\item First, in the first chapter of the book of Moses, Moses several times is ``in the similitude'' of Christ.
							\item We read in \hyperref[MOSES-1-6]{Moses 1:6}, ``thou art in the similitude of mine Only Begotten.''
							\item Then in \hyperref[MOSES-1-13]{Moses 1:13} Moses himself says, ``I am a son of God, in the similitude of his Only Begotten.''
							\item And finally, in \hyperref[MOSES-1-16]{Moses 1:16}, Moses again says, ``Get thee hence, Satan; deceive me not; for God said unto me: Thou art after the similitude of mine Only Begotten.''
							\item So here Moses, a human being, is identified as being ``in the similitude'' of Christ, and the reference here seems to be to the fact that Moses is a child of God in some sense the same way that Christ is.
							\item ``Similitude'' doesn't mean exactness, but there's some kind of similarity.
							\item The idea here might stem from the comment made in \hyperref[GENESIS-1-26]{Genesis 1:26}, about the creation, where we read, ``And God said, Let us make man in our image, after our likeness.''
							\item (In the Hebrew of that verse, ``image'' is \he{.sElEM}, and ``likeness'' is \he{d*:mUt}. \he{d*:mUt} is translated as ``similitude'' twice in the KJV, at \hyperref[2CHRONICLES-4-3]{2 Chronicles 4:3} and \hyperref[DANIEL-10-16]{Daniel 10:16}, so there is a linguistic connection.)
							\item (In the LXX of the Genesis verse, ``similitude'' is \gr{ὁμοίωσις}. The only appearance of that word in the NT is at \hyperref[JAMES-3-9]{James 3:9}, where it is ``similitude'' in the KJV. The James verse looks like a reference to the verse in Genesis.)
							\item The next verse that can help us with this passage in 138 about ``sacrifice'' is in the Book of Mormon, at \hyperref[JACOB-4-5]{Jacob 4:5}. Here we learn about the prophets before Nephi's people, and their attitude toward Christ. Let's read from verses 4--5.	
							\item This verse is directly about sacrifice---specifically, it's about Abraham's sacrifice.
							\item \tr{What, exactly, is the ``similitude of God and his Only Begotten Son,'' as described in this verse?}
							\item Finally, we go back to the Book of Moses for even more help on this idea of people offering a sacrfice ``in the similitude'' of Christ's great sacrifice.
							\item We need to read \hyperref[MOSES-5-5]{Moses 5:5--8}.
							\item (Discuss these verses.)
							\item \tr{What is it about this sacrifice that makes it ``in the similitude'' of Christ's sacrifice?}
						\end{compactitem}
					\item Let's continue to think about this sacrifice by looking at a few other verses.
						\begin{compactitem}
							\item While there are various ways of understanding what this sacrifice actually is, the scriptures suggest at least two pretty directly.
							\item The first is our baptism and the ``crucifixion'' of the old version of us, or the version of us that existed prior to being born again.
							\item This process is described in \hyperref[ROMANS-6-3]{Romans 6:3--6}.
							\item (Read and discuss those verses.)
							\item The other is the sacrifice that we are commanded to offer by Christ himself in the Book of Mormon.
							\item Let's read \hyperref[3NEPHI-9-19]{3 Nephi 9:19--20}.
							\item (Read and discuss those verses.)
						\end{compactitem}
					\item \tr{The verse seems to imply that you \textit{already have} a broken heart and a contrite spirit, and then you offer them \textit{as} the sacrifice. Is this right? What specific things do we \textit{do} to get those things?}
					\item \tr{How, specifically, do we make this sacrifice day by day? What are the specific things that we do and experience in order to offer this sacrifice?}
					\item In thinking about this commandment to make this sacrifice, I am struck by the words that McKay sees written in gold at the end of his vision.
					\item The words speak of the people who enter into the celestial kingdom, and they say, ``These Are They Who Have Overcome The World -- Who Have Truly Been Born Again.''
					\item We've just been talking about being born again; let's now think about the idea of ``overcoming the world.''
					\item \tr{As part of what we do with Christ in our salvation, what do you think it means to ``overcome the world''?}
					\item In our ``overcoming'' of the world, we follow the example of Christ.
						\begin{compactitem}
							\item At \hyperref[JOHN-16-33]{John 16:33}, Christ says, ``In the world ye shall have tribulation: but be of good cheer; I have overcome the world.''
							\item \hyperref[1JOHN-5-4]{1 John 5:4--5} then connects the ideas that we've been talking about, to tell us about how \textit{we} overcome the world.
							\item \tr{How do you continue to make this overcoming real in your own life?}
						\end{compactitem}
					\item \textbf{Other questions to ask or topics to discuss:}
						\begin{compactitem}
							\item \tr{What does it mean that the Lord ``will judge all men according to their works, according to the desire of their hearts''? Is this a good thing for us, or bad?}
							\item \tr{And is it according to \textit{both} our words \textit{and the desires of our hearts? Or what is the relationship between those things?}}
						\end{compactitem}
					
					
					
					\item Note that McKay's vision ends the semi-circle words in gold, ``.''
					\item Note that one of the privileges of the Melchizedek priesthood is ``to commune with the general assembly and church of the Firstborn.''
					\item The biblical source for this language is \hyperref[HEBREWS-12-23]{Hebrews 12:23}, but there the term ``firstborn'' is plural, not singular.
					\item (\hyperref[HEBREWS-1-6]{Hebrews 1:6} speaks of Christ as the ``firstbegotten,'' and you see an important connection here to \hyperref[EXODUS-4-22]{Exodus 4:22--23}. The fact that the followers of Christ are then called the ``assembly of the firstborn'' shows that they have a connection to Christ, or a certain relationship with him; that is the new birth that McKay is talking about in his vision.)
					\item (There are three uses of \gr{πρωτοτόκος} in Hebrews: 1:6, \hyperref[HEBREWS-11-28]{11:28}, and 12:23. The 11:28 is an interesting connection that forms a bridge from the Exodus to Hebrews and then to 12:23.)
					\item Connect to DC 138 by means of verse 13; see all the notes there, including stuff on the broken heart and contrite spirit; pay close attention to the Romans verses.
					\item (I think I also want to connect to ``overcoming the world,'' which is in the McKay vision at the end)
					\item Note that we already have to HAVE a broken heart and a contrite spirit in order to offer those things as the sacrifice.
					\item I think different things are spoken of as this sacrifice in different ways, but much of it has to do with the new birth.
					\item NOTES FROM AB (Koester) Hebrews 12:23, page 545: ``12:23. and the assembly. This word (ekklésia) was used for Israelite gatherings (Deut 31:30; Judg 20:2), Christian congregations (e.g., | Cor 1:2; Gal 1:2; Rev 1:4), and the whole church (Matt 16:18; Acts 9:31). In classical Greece the assembly had power to make decisions for the city (OCD, “ekklésia”; MM, 195). In the Greco-Roman world some decisions were made by a municipal council (boulé) and others by the “assembly” of adult male citizens (Acts 19:39). In Hebrews all the righteous, from Abel on, belong to the heavenly assembly (Heb 12-236); of the firstborn. The firstborn had special rights of inheritance and blessing, which Esau lost (NOTE on 12:16), but which the firstborn in heaven receive. The whole people of God are God’s firstborn (Exod 4:22—23). Although the firstborn in Israel were spared death in Moses’ time (Heb 11:28), those who belong to God’s city do die (cf. 11:13, 21-22, 35-37), but are delivered through resurrection, as was Jesus, who is God’s “firstborn” in a singular sense (1:6).''
					\item Firstborn (in Hebrew), same root as ``firstfruits''; many connections. 
					\item We have to offer broken heart/contrite spirit because these are things we \textit{must offer}; they cannot be taken from us by anyone else.
					\item Being judged ``according to their works according to the desires of their hearts''
						\begin{compactitem}
							\item Other verses that are relevant here?
							\item \hyperref[1PETER-4-6]{1 Peter 4:6} could be relevant here.
						\end{compactitem}
					\item Thrice, ``difference between sleight of hand and giving everything you have''
				\end{compactitem}
				
			\item Stuff about celestial kingdom, including church of the firstborn? DC 76 stuff? 
				\begin{compactitem}
					\item ``Degrees'' of the celestial kingdom?
				\end{compactitem}
				
			\item Fold the stuff from 138 into these things about DC 137?	

			\item Spirit world/preaching to the dead
				\begin{compactitem}
					\item Note the AB volume argues that this cannot refer to Christ preaching to the spirits of deceased humans. 
				\end{compactitem}
		
		\end{compactitem} 
	
	% ----------------------
	\subsection{2 November 2025: Sunday school on DC 124}
	% ----------------------
		\label{TALKS-02-11-2025}
		
		\begin{compactitem}

			\item Hi everyone.
				\begin{compactitem}
					\item Today we're talking about DC 124.
					\item We're going to begin by getting a little background.
					\item Then I would like to talk about the idea of a ``stake''---what is a \textit{stake}, exactly, and why is this imporant?
					\item Then I want to consider some connections to the temple.
					\item That's the plan for today.
				\end{compactitem}
				
			\item Background to section 124
				\begin{compactitem}
					\item On 19 January 1841, Joseph Smith dictated a revelation in Nauvoo, Illinois, designating the city as the new gathering place for the Latter-day Saints. 
					\item The revelation began by addressing Joseph Smith personally before expanding to instruct other individuals and the Saints in general. 
					\item Over time, the lengthy revelation served as a sacred charter for the Saints in Nauvoo, paralleling the recently passed act of incorporation that functioned as the city’s secular charter. 
					\item While the city charter represented the church’s legal permanence in Nauvoo, the revelation reassured members that the move to Illinois did not mean abandoning their mission to build Zion in Missouri. 
					\item The revelation declared that Nauvoo would be ``a corner stone of Zion,'' or a ``stake'' of Zion, rather than Zion itself. 
					\item The city was not to be a temporary refuge, as shown by the commandment to build a temple there. 
					\item The directive to build a temple in Nauvoo gave divine authority to instructions Joseph Smith had already been giving publicly for over six months. 
					\item The revelation emphasized the importance of the temple by stating that certain ordinances, such as baptisms for the dead, could only be performed within it. 
					\item An October 1839 church conference had already identified Nauvoo as a new gathering place for the Saints, and the revelation reaffirmed this role by calling it the cornerstone of Zion. 
				\end{compactitem}
				
			\item The theology of \textit{stakes}
				\begin{compactitem}
					\item Let's read \hyperref[DC124-2]{verse 2}.
						\begin{compactitem}
							\item We read that this ``stake,'' apparently meaning the stake of Nauvoo, has been ``planted'' to be a ``cornerstone of Zion.'' 
						\end{compactitem}
					\item First, why do we call these geographical boundaries ``stakes''?
					\item Where does this language come from?
						\begin{compactitem}
							\item (See below for Isaiah 54:2 and Isaiah 33:20.)
							\item For the Hebrew, see \hyperref[STAKE]{Stake}.
							\item But the word is used in Exodus for the ``pins'' or ``pegs'' that hold the tabernacle together.
							\item See for example \hyperref[EXODUS-27-19]{Exodus 27:19}.
							\item \hyperref[EZRA-9-8]{Ezra 9:8}, the temple language is all over this.
						\end{compactitem}
					\item Strange that a stake of a tent is a cornerstone of something?
					\item There may be a parallel here to what happens in the OT---David takes Jerusalem, and brings the ark in (2 Samuel 6), and it dwells in the ``curtains,'' which are still like the tent (\hyperref[2SAMUEL-7-2]{2 Samuel 7:2}).
					\item Note the parallel there to \hyperref[DC101-21]{DC 101:21} (that's December 1833).
					\item But then his son Solomon builds the actual temple, so there is a transition here from the stakes of the old tent or ``tabernacle,'' to a permanent structure.
					\item The imagery of stakes suggests wooden stakes planted in the ground throughout many different places, which would then allow you to secure a tent or covering over something; the curtains in 101:21 are part of that.
					\item See the temple as the culmination of the theology of stakes. It's the place where God visits, just like the ancient tent or tabernacle.
					\item And what we see in the history of the restoration is a transition from the ``stakes of Zion,'' or these more ``temporary'' structures, to the permanent ones: first in Kirtland, and then planned in Jackson County in Missouri but never actually built, and then again commanded to be built in Nauvoo.
				\end{compactitem}
				
			\item Extra notes on ``stake''
				\begin{compactitem}
					\item The word ``stake'' only appears twice in the OT, both plural, at \hyperref[ISAIAH-33-20]{Isaiah 33:20} and \hyperref[ISAIAH-54-2]{Isaiah 54:2}. At 33:20 the Hebrew is \he{yAted}, and at 54:2 it's the same; see \hyperref[YATED]{\hel{yAted}}.
					\item Isaiah 22:23: ``And I will fasten him as a NAIL in a sure place; and he shall be for a glorious throne to his father’s house.''
						Zechariah 10:4: ``Out of him came forth the corner, out of him the NAIL, out of him the battle bow, out of him every oppressor together.''
						Ezra 9:8: ``And now for a little space grace hath been shewed from the LORD our God, to leave us a remnant to escape, and to give us a NAIL in his holy place, that our God may lighten our eyes, and give us a little reviving in our bondage.''
					\item ``weight of glory,'' for \hyperref[ISAIAH-22-23]{Isaiah 22:23}?
					\item The word ``stake'' appears twice in the BOM, and both are references to Isaiah 54:2.
					\item ``Stake'' does not appear at all in the PGP.
					\item ``Stake'' appears 25 times in the DC.
						\begin{compactitem}
							\item The first is two occurrences in 68 (25 and 26), which is 1 November 1831. \textbf{HOWEVER}, both of these occurrences are later additions and are not in the originals.
							\item It then has two occurrences in 82 (13 and 14), which is 26 April 1832.
							\item The next is 94:1, 2 August 1833.
							\item Note, though, that the term also appears in DC 133, which according to the JSP is 3 November 1831. So there is a source going back to those first few days in November 1831 for ``stakes.'' Remember, see above, the uses in 68 are later additions.
							\item By mid-1833, in a letter to church leaders, JS and the first presidency can say, ``Kirtland, the Stake of Zion, is strengthening continually'' (\hyperlink{EMS-1833-JS-25-JUN-1833-e}{see here}). 94:1 then also says ``stake of Zion.''
							\item \hyperref[DC101-21]{DC 101:21} might be the most helpful verse of all. The JSP version says, ``behold there [is] none other place appointed than that which I have appointed neither shall there be any other place appointd that [than] that which I have appointed for the work of the gathering of my saints until the day cometh when there is found no more room for them and they shall be called stakes for the curtains of Zion or strength of Zion.''
							\item The ``curtains'' language also seems to come from Isaiah 54:2. There it's \he{y:riy`Ah}; HALOT 439a/PDF 495 ``tent curtain, tent (for the ark).'' The Isaiah verse is listed in the first meaning, the second one applies to 2 Samuel 7:2.
							\item The ``curtains'' are supposed to be ``a covering upon the tabernacle'' (Exodus 26:7).
							\item (need more about curtains next)
						\end{compactitem}
				\end{compactitem}
				
			% Notes
				% STAKE
					% YATED
						% NIDOTTE 2:568
						% TWOT 418b
				
			\item But things are progressing so quickly that the stakes can't become cornerstones fast enough, and so this temple work and these temple experiences happen outside of the finished building, in structures that are sort of temporary temples, just like the tent of the congregation or tabernacle was a temporary structure while they didn't have a permanent one.
				\begin{compactitem}
					\item Let's read \hyperref[DC124-29]{DC 124:29--30}.
					\item So what pattern do we see here?
					\item Let's also now read \hyperref[DC124-31]{DC 124:31--33}.
					\item So this was a temporary thing.
					\item We see another example of this in the original presentation of the endowment in May 1842. Does anyone know where it happens?
					\item Joseph Smith journal for May 4, 1842: ``In council in the Presidents \& General offices with Judge [James] Adams. Hyram [Hyrum] Smith Newel K. Whitney. William Marks, Wm Law. George Miller. Brigham Young. Heber C. Kimball \& Willard Richards. [illegible] \& giving certain instructions concerning the priesthood. [illegible] \&c on the Aronic Priesthood to the first [illegible] continueing through the day.''
					\item This short entry then gets expanded into a longer comment by Willard Richards for Joseph Smith's history:
						\begin{compactitem}
							\item ``JS instructed those present “in the principles and order of the priesthood, attending to washings \& anointings, endowments, and the communications of keys, pertaining to the Aronic Priesthood, and so on to the highe[s]t order of the Melchisedec Pristhood, setting forth the order pertaining to the Ancient of days \& all those plans \& principles by which any one is enabled to secure the fulness of those blessings which has been prepared for the chu[r]ch of the first-born, and come up, and abide in the prese[n]ce of Eloheim in the eternal worlds. In this council was institutd the Ancient order of things for the fir[s]t time in these last days....'' These instructions ``were of things spiritul, and to be received only by the spiritual minded: and there was nothing made known to these men but what will be made known to all saints, of the last days, so soon as they are prepared to recive, and a proper place is prepared to communicate them, even to the weakest of the saints; therefore let the saints be diligent in building the Temple and all houses which they have been or shall hereafter be commanded of God to build, and wait their time with patience, in all meekness faith, \& perseverance unto the end, knowing assuredly that all these things refer[re]d to in this council are always governd by the principles of Revelation.''
						\end{compactitem}
					\item I really want to emphasize the comment made here that says that part of the blessings of the endowment is to ``abide in the prese[n]ce of Eloheim in the eternal worlds.''
					\item The purpose of the temple is to bring is into God's presence.
						\begin{compactitem}
							\item This is sometimes hard to keep in mind because we talk about temple ``work,'' and going to the temple \textit{always} involves officiating in the ordinances for someone else.
							\item But don't let this get in the way of you understanding what has always been true about the tabernacle, and Solomon's temple, and the Kirtland temple, and Joseph Smith's store, and the Nauvoo temple, and every temple now: it is about preparing a place where God can dwell on earth.
							\item We read this in \hyperref[DC124-27]{DC 124:27}.
						\end{compactitem}
				\end{compactitem}
				
			\item Why do we have this house?
				\begin{compactitem}
					\item \hyperref[DC124-40]{DC 124:40}; to reveal ordinances.
					\item \hyperref[DC84-20]{DC 84:20--22}, about ordinances being necessary to experience God's presence.
					\item \textit{WHY} are ordinances necessary for us to ``see the face of God, and live''?
					\item They are necessary \textit{for us}, because human beings \textit{need reminders} and \textit{images} and \textit{metaphors} and \textit{representations} in order to be able to understand and do things.
					\item Sanctification occurs by powers that we cannot control; Christ's power.
					\item The ordinances and rituals are an aid to us in this sanctification process.
					\item 
					\item If you don't do these things, ``I will not perform the oath''---\hyperref[DC124-47]{DC 124:47}.
						\begin{compactitem}
							\item ``Oath,'' \hyperref[DC84-39]{DC 84:39}. 
							\item God will unveil his face to you in his own time, \hyperref[DC88-68]{DC 88:68}.
							\item Endowment ceremony at the veil, this person ``having conversed with the Lord through the veil, desires now to enter his presence. Let him enter.''
						\end{compactitem}
				\end{compactitem}
				
			\item What is the ``fulness of the priesthood'' (\hyperref[DC124-28]{DC 124:28})?
				\begin{compactitem}
					\item The second anointing, in fall of 1843.
					\item We are anointed in the endowment only to \textit{become} kings and priests, or queens and priestesses.
					\item There will apparently be a moment when we are anointed to actually \textit{be} those things.
				\end{compactitem}
				
			\item DC 121
				\begin{compactitem}
					\item Starting at \hyperref[DC121-35]{DC 121:35}?
					\item Verses \hyperref[DC121-45]{45--46} in particular.
				\end{compactitem}
		
		\end{compactitem}
	
	% ----------------------
	\subsection{21 September 2025: Sunday school on DC 103--105}
	% ----------------------
		\label{TALKS-21-09-2025}
		
		\begin{compactitem}
		
			\item Hi everyone.
				\begin{compactitem}
					\item Welcome to class today. I'm looking forward to talking about the scriptures.
					\item I'm really happy that Sarah is teaching, and I was honestly pretty devastated two weeks ago when I had to miss the discussion on DC 93.
					\item That section is one of my favorite passages of all scripture, and one of the deepest and richest, and I have loved our discussions in this group and so I was very sad to miss it.
					\item But I know me, and I'm sure I'll find a way to connect our future discussions back to section 93 somehow.
				\end{compactitem}
				
			\item In fact, that might happen today---our topic is DC 103--105, and we'll be looking at some passages from these sections as well as making connections to some other places in the scriptures that can help us think about how the gospel is operating in these events.
				\begin{compactitem}
					\item Section 103 is a revelation received on February 24, 1834. 
					\item The same day that this revelation was dictated, Parley P. Pratt and Lyman Wight, who had traveled from Missouri to Kirtland, Ohio, reported to the Kirtland high council on the condition of the Missouri members, most of whom had taken refuge in Clay County.
					\item Pratt and Wight also asked the high council ‘how and by what means Zion was to be redeemed from our enemies.’
					\item JS then volunteered to lead an expedition to Missouri to assist in Zion’s redemption, and thirty or forty others stated they would go with him.
					\item Before adjourning, the council appointed JS as ‘Commander in Chief of the Armies of Israel and the leader of those who volunteered to go.’
					\item ``Extant records do not state whether this revelation was dictated before, during, or after the high council meeting.
					\item It is difficult to determine whether the revelation motivated the decisions of JS and others in the high council or whether those decisions were affirmed by the revelation’s dictation.
					\item ``JS and others quickly acted on the revelation’s directions.
					\item Just two days later, he and Parley P. Pratt ‘started from home to obtain volenteers for Zion,’ and Orson Pratt and Orson Hyde did the same.
					\item Lyman Wight and Sidney Rigdon, also appointed in the revelation as recruiters, apparently left Kirtland on 28 February.
					\item It is not clear when Hyrum Smith and Frederick G. Williams, the other pair assigned in the revelation to recruit members for the expedition, departed, but John Murdock reported that Hyrum Smith was at Edmund Bosley’s home in Livonia, New York, on 15 March 1834.
					\item For the next several weeks, these eight men held meetings in which they preached, recruited volunteers, and raised funds to help restore the refugees to their homes in Jackson County.
					\item These activities ultimately resulted in the formation of the Camp of Israel, an expedition of more than two hundred individuals that marched to Clay County, Missouri, in the summer of 1834. 
				\end{compactitem}
				
			\item So all of a sudden they're just taking off, and kind of running around with this militia thing, and when I read about this the first thing I think is, who is taking care of the kids?
				\begin{compactitem}
					\item I'm really sensitive to this, because taking care of our kids is so much work and it's so tiring.
					\item I kind of imagine Joseph and Hyrum Smith coming home to a bunch of whiny, difficult kids and telling Emma and Jerusha, ``We're going bowling.''
					\item And the wives are like, ``Please get these kids away from me before I have to bury them out back.''
					\item And the husbands say something like, ``But we \textit{always} go bowling on Wednesday nights.''
					\item Except in this case, their bowling league is actually hundreds of miles away, and they have to walk there, and they have to leave right now to gather up all the bowling team members into a group.
					\item This is not just a point about paying tribute to childcare and other domestic labor that made many of these early events in the restoration possible, although it is that.
					\item It's also a larger point about how the material conditions of a culture and a society affect the way the restoration unfolds, and continues to unfold even today.
					\item When we talk about the ``Camp of Israel,'' or all this ``Zion's Camp'' stuff, we almost always point out that a huge number of early church leaders were among that group that marched. 
					\item But we also have to remember that in order for \textit{them} to march, a bunch of other people \textit{couldn't} march, and so of necessity, a bunch of people could not have been involved in this experience that we say was so important in forging early leadership.
					\item When we think about why things are the way they are in the church, we have to always keep in mind how the material conditions of the restoration demanded that certain cultural practices become a part of church organization---in this case sex-segregated labor roles.
					\item Almost two hundred years later, the material conditions in many ways are now very different, and so when we ask questions about why things are set up the way they are, we tend to look back in the history and look for deliberate reasons or purposes for why things \textit{should} be this way.
					\item But in reality, a lot of the time the answer is not that intentional, since what's really happening is that people in the past were playing out the scripts given to them by those material conditions of culture and society.
					\item And the same is true now.
				\end{compactitem}
				
			\item Now let's get into the scriptures. We are going to begin reading in section 103.
				\begin{compactitem}
					\item We'll start with \hyperref[DC103-3]{103:3--4}.
						\begin{compactitem}
							\item According to these verses, the ``chastening'' and ``chastisement'' the people were experiencing were due to them not ``hearkening'' to the precepts and commandments they had received.
							\item \hl{Is this a general principle that's for the most part true in our lives as well? Is it true in general that our mistakes or lack of obedience are responsible for the difficult things we experience?}
							\item \hl{Or is this just a comment that's limited to these particular circumstances?}
						\end{compactitem}
					\item The revelation makes a similar remark in \hyperref[DC103-7]{verses 7--8}. If the people don't keep all the commandments and don't ``hearken to observe'' all the words, then the ``kingdoms of the world shall prevail against them.''
					\item \hyperref[DC103-9]{Verses 9--10} also are about roughly the same idea, and reference Christ's words in the Sermon on the Mount at \hyperref[MATTHEW-5-13]{Matthew 5:13}. 
						\begin{compactitem}
							\item There's an interesting pun in these verses, with the words ``savior'' and ``savor,'' and of course this pun only works in English; it doesn't work in the original Greek of the New Testament (``saviour'' = \gr{σωτῆρος}, ``savor'' = aorist passive of \gr{μωραίνω}).
							\item The Lord's people are referred to as ``saviors of men'' here. Verse 9 says that these people were ``set to be'' the ``saviors of men.''
							\item \hl{How is it possible for us to be ``saviors'' of other people?}
							\item This is interesting, as apart from these two verses here, there is only one place in the DC where the word ``savior'' is applied to someone other than Jesus Christ.
							\item That is at \hyperref[DC86-11]{DC 86:11}. These verses in 103 seem clearly a reference to that verse in section 86, as that verse also talks about being a ``light unto the Gentiles,'' and ``light'' is spoken of here in 103 too.
						\item In 86:11 we learn that it is ``through this priesthood'' that people are these ``saviors,'' and so we play a role in the salvation of other people when we help them access more of the power of Christ's atonement through ordinances of the priesthood.
						\end{compactitem}
				\end{compactitem}
				
			\item These verses are preliminary to the larger purpose of the revelation, which is to talk about people getting thrown out of ``Zion,'' or certain areas of Missouri.
				\begin{compactitem}
					\item There's more Biblical language to describe these events.
						\begin{compactitem}
							\item \item \hyperref[DC103-13]{Verse 13} talks about the people's ``restoration to the land of Zion,'' and says that they are ``no more to be thrown down.''
							\item The next verse, \hyperref[DC103-14]{14}, also uses the idea of being ``thrown down,'' and this phrase is common in the Old Testament, especially in talking about ``altars,'' or sources of religious authority, that are destroyed or toppled.
							\item Here the reference may also be to the parable of the nobleman and the watchmen in \hyperref[DC101]{DC 101} (especially verses 57 and 75).
						\end{compactitem}
					\item In any case, we read in \hyperref[DC103-15]{verse 15} that ``the redemption of Zion must needs come by power.''
					\item \hl{Could somebody please read verses 15--18 for us?} (\hyperref[DC103-15]{DC 103:15--18})
						\begin{compactitem}
							\item God is going to ``raise up'' a person, presumably Joseph Smith in his capacity as ``Commander in Chief,'' to lead them ``as Moses led the children of Israel.''
							\item Now here, the revelation begins to use language not just from the Old Testament in general, but specifically from the book of Exodus, to describe the situation of these early saints, and to describe how the Lord will interact with them in order to begin resolving this immense challenge that they're facing.
							\item So for example, let's look at \hyperref[DC103-17]{verse 17} in particular.
							\item The revelation says that these early saints ``are the children of Israel.''
							\item It says that the people are ``of the seed of Abraham,'' and so makes a connection to the promises made to the ``seed of Abraham''; this also happens in \hyperref[EXODUS-33-1]{Exodus 33:1}: Moses brought the children of Israel ``out of the land of Egypt, unto the land which I sware unto Abraham, to Isaac, and to Jacob, saying, Unto thy seed will I give it.''
							\item We also read in 103:17 that the people will be ``led out of bondage by power, and with a stretched-out arm.''
							\item Similarly, just prior to actually leaving Egypt, at \hyperref[EXODUS-13-3]{Exodus 13:3} Moses says, ``Remember this day, in which ye came out from Egypt, out of the house of bondage; for by strength of hand the LORD brought you out from this place.''
							\item Also \hyperref[EXODUS-13-14]{Exodus 13:14}: ``And it shall be when thy son asketh thee in time to come, saying, What is this? that thou shalt say unto him, By strength of hand the LORD brought us out from Egypt, from the house of bondage.''
							\item So just in verse 17 alone, we see multiple explicit references to the Exodus story of Moses leading the children of Israel out of Egypt.
							\item But we also know that this ``redemption of Zion'' did not really work out the way that Joseph Smith and others were hoping; they got out there in the camp and they were outnumbered by other militias, a bunch of people got cholera, and they had to basically turn around and go home without accomplishing the thing tehy were supposed to do.
							\item The story we tell about this is the one I mentioned at the beginning, about how, well, the expedition didn't really accomplish its explicit purposes, \textit{but} it did bond the people together, clarified who the next generation of leaders would be, and so on.
							\item I don't necessarily have a problem with viewing things this way, it just seems like a bit of after-the-fact revisionist history---like the expedition failed, and so we still need a reason to justify it, which means we look at the leadership point, even though the chances were that all of those men were going to be involved in early leadership anyway, regardless of whether this expedition had even happened or not.
							\item The question I have for you is, \hl{what is the lesson that *WE* are supposed to learn from Zion's Camp? Is there a principle here about human action and the Lord's purposes? Or is this just an isolated event in the history of the church?}
						\end{compactitem}
					\item Now let's read \hyperref[DC103-19]{DC verses 19--20}, because these are going to introduce a crucial point that not only connects us back to the story in Exodus, but also tells us something about why the restoration in our time is so important.
						\begin{compactitem}
							\item The Lord tells the people receiving the revelation to ``let not your hearts faint.''
							\item And then he makes kind of an unusual comment: he says, ``I am \textit{not} telling you something that I told people in the past,'' which was ``Mine angel shall go up before you, but not my presence.''
							\item He says, I'm \textit{not} telling you that---I told it to pepole in the past, but I'm not telling it to you.
							\item In the next verse, verse 20, he then clarifies what he \textit{is} telling to people \textit{now}: ``Mine angels shall go up before you, and also my presence, and in time ye shall possess the goodly land.''
							\item So notice the difference between what was said before, and what is being said now: The ``angel'' or ``angels'' were always going to ``go up before'' the people, but the ``presence'' was not going to before, whereas it is now.
							\item \hl{What would the difference be between the people just having the ``angels'' with them, and having God's presence with them? Do we experience a similar difference in our lives?}
							\item Here we also find some important connections to Exodus.
								\begin{compactitem}
									\item I'm going to read from \hyperref[EXODUS-32-34]{Exodus 32:34} first.
									\item Then \hyperref[EXODUS-33-1]{Exodus 33:1--3}.
									\item So here we find God saying something along the lines of what the revelation in DC 103 says---it will be the angels, not the presence of the Lord, that goes with them.
									\item And in fact the children of Israel had already requested something like this, earlier on.
									\item In \hyperref[EXODUS-20-18]{Exodus 20:18--21} they make a request of this kind. Let's read those verses.
									\item There the people experience certain manifestifations of natural phenomena, that are connected with God's presence.
									\item And they're like, ``No thank you.''
									\item But Moses forges ahead into the darkness, and talks with God.
									\item \hyperref[DEUTERONOMY-5-23]{Deuteronomy 5:23--27} explains more of the reasons---they were afraid of \textit{dying}.
									\item Back in Exodus 33, we still see these different attitudes toward God's presence.
									\item Moses sets up the tabernacle, and the people go out to it, but when Moses goes out to it, that's when the ``cloudy pillar'' descends.
									\item That's when we read things like \hyperref[EXODUS-33-9]{Exodus 33:9 and 11}, where Moses is speaking to God ``face to face.''
									\item This is a tricky question that I hope makes sense: \hl{do we find these parallel attitudes in ourselves, or in our times? The attitude of saying, ``No, I'd like someone else to be an intermediary to God for me, please,'' and the attitude of, ``I'm going to walk into the darkness to experience God's presence?''}
								\end{compactitem}
							\item What impresses me about the revelation in section 103 is that this talk about God's ``presence'' is foreshadowing something much more important than Zion's Camp. It's foreshadowing the temple, as the place where we go to experience God's presence, and foreshadowing the endowment, or the gift of power and knowledge that we receive by coming into that presence.
								\begin{compactitem}
									\item \hyperref[DC105-9]{DC 105:9--12}
									\item \hyperref[DC105-18]{DC 105:18--19}
									\item \hyperref[DC105-33]{DC 105:33}
									\item The whole purpose of the restoration is to get \textit{individual people} into God's presence, \textit{without} a mediator. It is to get \textit{you} to experience God in this direct way, and not just be on the sidelines, like the children of Israel are when they watch Moses enter the tabernacle and the cloudy pillar comes down.
									\item \hyperlink{EMS-1839-JS-CIR-26-JUN-CIR-2-JUL-1839-WW-c}{Joseph Smith on people experiencing what he experienced}: ``God hath not revealed any thing to Joseph but what he will make known unto the Twelve \& even the least saint may know all things as fast as he is able to hear them for the day must come when no man need say to his neighbour know ye the Son for all shall know him (who remain) from the least to the greatest''
									\item (stuff about how ``presence'' is ``face'' in Hebrew)
									\item Moses with endowment-like presence experience, and shining face: \hyperref[EXODUS-34-29]{Exodus 34:29--35}.
									\item See God's face: \hyperref[DC93-1]{DC 93:1}, \hyperref[DC93-19]{DC 93:19}, 
									\item We face a similar choice: do we want to enter God's presence ourselves---do we want to face down the fear of things outside of us and inside of us, and experience God's presence more directly, without mediation?
									\item (The people feared death, and \hyperref[EXODUS-33-20]{Exodus 33:20} says that Moses can't see God's face, because ``there shall no man see me, and live.'' But this is a problem explicitly addressed in the restoration.)
									\item (\hyperref[DC84-21]{DC 84:21--22}: ``And without the ordinances thereof, and the authority of the priesthood, the power of godliness is not manifest unto men in the flesh; For without this no man can see the face of God, even the Father, and live.'')
									\item (\hyperref[DC88-67]{DC 88:67--68}: ``And if your eye be single to my glory, your whole bodies shall be filled with light, and there shall be no darkness in you; and that body which is filled with light comprehendeth all things. Therefore, sanctify yourselves that your minds become single to God, and the days will come that you shall see him; for he will unveil his face unto you, and it shall be in his own time, and in his own way, and according to his own will.'')
									\item (\hyperref[DC76-113]{DC 76:113--118})
								\end{compactitem}
						\end{compactitem}
				\end{compactitem}
		
		\end{compactitem}
	
	% ----------------------
	\subsection{16 March 2025: Sacrament meeting talk on Christ as the source}
	% ----------------------
		\label{TALKS-16-03-2025}
		
		Alicia didn't speak with me this time, it was just me for this one.
		
		\begin{compactitem}
		
			\item Hi everyone. 
				\begin{compactitem} 
					\item It's finally warming up. I have been really happy about this because the winter was difficult, with the kids being out of school so much. 
					\item We did a lot of sledding and drank a lot of hot chocolate but actually we had a pretty hard time.
					\item I don't know how y'all have been doing, but in reality the hard time started before all the snow.
					\item I at least feel like I have been barely functional for months.
					\item I feel like I'm a zombie wandering from thing to thing, doing the bare minimum of whatever it is I need to do in order for things to turn out correctly.
					\item And do you know what has done the most for me during this time? Do you know what has granted me the greatest amount of strength and endurance and will to survive?
					\item Mountain Dew.
					\item Seriously, I am alive today because of Mountain Dew.
					\item Not in the sense that Mountain Dew keeps me alive, but in the sense that it makes my life worth living.
					\item You know how the scripture in Moses 1:39 says, ``For behold, this is my work and my glory---to bring to pass the immortality and eternal life of man.''
					\item If a similar statement were written about me, it would say, ``For behold, this is my work and my glory---to bring to pass the acquisition and consumption of Baja Blast.''
				\end{compactitem}
				
			\item I've been thinking a lot about why I feel barely functional, because it hasn't always been this way, or when I've had stretches like this in the past, they haven't always lasted this long.
				\begin{compactitem}
					\item (Claude the AI playing Pokemon? Doing the same things over and over again?)
					\item I think a big part of it is that, despite all the Mountain Dew---or maybe because of it---I just feel like I have so little energy.
					\item I'm just so tired all the time, it feels like.
					\item Unfortunately, sometimes my tendency is to handle this fatigue by eating.
						\begin{compactitem}
							\item Some days it's like I'm living out a Jim Gaffigan bit in real life.
							\item ``Gee, all this eating has made me really tired...welp, I guess I better eat.''
						\end{compactitem}
					\item But it isn't just the food. 
					\item And right now, for whatever reason, it feels like there are a lot of things in my life that just drain all that energy away without replacing it.
					\item There are a lot of things that are just energy sinks, in that I'm pouring energy into them, but sometimes not getting a lot back.
					\item So whatever energy I have, gets siphoned away here, gets siphoned away there, and before you know it it's 2:00pm and I've already been zombified and I'm just like, how can I spend the rest of the day just holding things together? 
					\item And in the midst of all this I keep thinking, I should be able to do more than this. I used to do more than this; I used to be more than barely functional. I have all these goals and plans, and they're good things, and yet it feels like I just can't bring myself to do them. 
				\end{compactitem}
				
			\item There's an idea in physics that's related to all this, and that is the idea of \textit{entropy}.
				\begin{compactitem}
					\item Entropy has a lot of technical definitions in different areas of physics, but the basic idea is that entropy is a measure of how spread out energy is.
					\item In order to accomplish something meaningful, you don't want your energy spread out all over the place. You want it concentrated and organized into a form that is useful for what you want to accomplish, like when you have the gas in your car all in the tank. That's where you want it, because it's concentrated there.
					\item When you add energy into something, like a machine or something like that, that energy you add produces something---it makes something happen.
					\item But the laws of physics and this idea called entropy have some bad news for us.
					\item Those things guarantee that you will \textit{never} get \textit{more} energy out of something than you put into it. That's impossible.
					\item But it's actually worse than that, because you also can't even get the \textit{same} amount of energy out of something that you put into it: you \textit{always} lose some of it.
					\item When you put energy into a machine, not all of it gets turned into useful output---some is always lost as waste heat, or as friction, or just as inefficiencies in the way the machine works.
					\item Over time, the energy that you add spreads out, it becomes less organized and less concentrated, and so as far as you're concerned it's just lost to you. It's just gone and you can't get it back.
					\item Now I know that when we talk about ourselves and the ways we live and our own experiences, this idea from physics doesn't necessarily apply in the exact same way.
					\item But this image has really been impacting me lately: I put all this energy that I have into all these different things, and it just feels like I sometimes get so little of it back.
					\item It's like I'm just dumping it into all the ``machines'' of my life, and then it gets used a bit and then spread out and unconcentrated, disorganized, and it's all lost to me.
					\item And I'm left wondering, what happened here? Am I doing something wrong? Is this just the way it is? 
				\end{compactitem}
				
			\item Of course, not everything is like this. Not everything is just an energy sink in this way, where it feels like 
				\begin{compactitem}
					\item There are probably things that do energize you---they fill you up rather than drain you.
					\item And these things will be different for different people.
					\item They aren't always necessarily the things you want or expect, though.
					\item For example, I would really like for church on Sundays to be one of these things for me, but the truth is that it's often not.
					\item I would love to feel like I come here drained and leave here filled, but that usually isn't true.
					\item There are various reasons why this is, including the special experience of trying to help kids sit through sacrament meeting.
					\item It also has to do with how much I am able to actually mindfully prepare myself for coming, which again for me depends on what's going on at my house with all the kids.
					\item Whatever the things that give you more of a return on the energy you put in, hold on to those. And when they change, as they definitely have for me as time as marched on, I hope you're able to come across new ones.
				\end{compactitem}
				
			\item That isn't the message that I want to share with you today, though.
				\begin{compactitem}
					\item The message I want to share today is about how Jesus Christ fits into all of this.
					\item As I've been thinking about this and the very minimal level of function that I've been able to sustain in the past little while, some scriptures have come to my mind, about the Savior, and about the kind of role he's able to play in our life when we look for a certain kind of relationship with him.
					\item And Kate Murri talked about something like this last week as well, when she spoke about ``drawing the water from the wells of salvation.''
					\item So I'd like to share some of those scriptures with you. 
				\end{compactitem} 
				
			\item Let's have a look at some of these verses that can speak to some of the issues we're dealing with here.
				\begin{compactitem}
					\item First let's look at this idea of the ``well of salvation.''
					\item There are various places we could look for this, but the thing that comes first to my mind is John 4, when Christ is speaking to the woman at the actual well.
					\item Let's start at \hyperref[JOHN-4-10]{John 4:10} and read until the end of verse 14. 
						\begin{compactitem}
							\item I'm struck by the woman's reaction in verse 11. She doesn't ask what it is for water to be ``living''; maybe at that time the idea was already familiar to her.
							\item What she does is look around and notice, you don't have any tools to obtain this water. You've got not bucket, no rope, no nothing. So how are you going to get that water?
							\item There's this element of impossibility here. She just can't conceive of how this is going to be brought about.
							\item And in his response to her, Christ does not address this issue. He doesn't explain how it is that he's able to make this happen, or get this water. 
							\item Instead, he tells her about the water itself, in verses 13 and 14.
							\item And what we learn about the water is another seeming impossibility: this water is such that, if you drink it, you will never thirst again.
							\item This is something that is completely outside of our experience. I mean even if Mountain Dew somehow had a brand of bottle water---I assume it would be called ``Dew Drops'' or ``Hydro Dew'' or something---we could drink that water and we'd still get thirsty pretty soon after.
							\item But the water that Christ gives us is not like that; it's different from every other water we could drink. It somehow satisfies us in such a deep, lasting way that we don't ever thrist again. 
							\item So there are at least two impossibilities here, or two things that seem to go against the normal course of things that we experience:
							\item (1) Christ is able to get the water without the normal means, and 
							\item (2) The water fills us in a completely different way, that nothing else can do.
							\item Both of these things give me hope that working with Christ is different from how we work with other things in our lives. 
							\item We were talking before about how these laws of physics guarantee that you cannot get more out of a system or out of an interaction than you put into it.
							\item But Christ is suggesting something else here. He's hinting at an energy source that goes way beyond our normal understanding, and that might even change some of those rules we were considering earlier.
						\end{compactitem}
					\item There are other places in the scriptures that speak of similar things.
					\item For example, at \hyperref[2NEPHI-9-50]{2 Nephi 9:50--51} and \hyperref[2NEPHI-26-25]{2 Nephi 26:25}, we read about ``wine and milk'' and ``milk and honey.''
						\begin{compactitem}
							\item 9:50 says, ``Come, my brethren, every one that thirsteth, come ye to the waters. And he that hath no money, come, buy and eat, yea, come, buy wine and milk without money and without price.''
							\item And then 26:25 says, ``Come unto me, all ye ends of the earth; buy milk and honey without money and without price.''
							\item Now the water is mentioned in the first verse---we're suppoed to ``come to the waters.''
							\item But after that it's a bit of an unusual combination of things---\textit{milk} and \textit{honey}? It sounds like a little kid's lunch or something. Throw in some goldfish and you've got what my daughter eats every Saturday.
							\item Now obviously if these verses in the Book of Mormon, had been translated from the gold plates today, it's fair to say that it would talk about Mountain Dew Live Wire: ``The rocks themselves are soda fountains!!!'' 
							\item In the context of the scriptures, though, we find this idea of ``milk'' and ``honey'' as a pair appearing a bunch of times in the Old Testament, and it's almost always when they're speaking about the promised land.
							\item They say, it's a land \textit{flowing} with milk and honey. It's this absolute abundance of goodness, of things that are so sweet and precious to you, an unlimited amount of this goodness. 
							\item Jacob expands on this in 9:51: ``Wherefore do not spend money for that which is of no worth, nor your labor for that which cannot satisfy. Hearken diligently unto me and remember the words which I have spoken, and come unto the Holy One of Israel, and feast upon that which perisheth not, neither can be corrupted, and let your soul delight in fatness.''
							\item Once again this seems like an impossibility: won't the cows run out of milk? Won't the bees run out of honey, if we're eating it all the time? Won't we run out of steaks and wings and pulled pork and stuff?
							\item But it doesn't seem so. It's once again this energy source that goes beyond our understanding, that seems contradictory to what we normally experience.
							\item And yet this is the promise.							
						\end{compactitem}
					\item Christ's expansion of the Sermon on the Mount in 3 Nephi also has a thought to add here.
						\begin{compactitem}
							\item In the version of this speech that he gives near Jerusalem, he says, ``Blessed are they which do hunger and thirst after righteousness: for they shall be filled'' (\hyperref[MATTHEW-5-6]{Matthew 5:6}).
							\item And the verse he gives to the Nephites is very similar---he says, ``And blessed are all they which do hunger and thirst after righteousness, for they shall be filled \textit{with the Holy Ghost}'' (\hyperref[3NEPHI-12-6]{3 Nephi 12:6}).
							\item So in both remarks he says that we will be \textit{filled}, and he adds in the Book of Mormon account that this ``filling'' comes by means of the Holy Ghost.
							\item But for me the emphasis is on the \textit{filling}. You are going to be \textit{filled}---you will not be left empty; you will not be left drained; you will not be left feeling like you are pouring your energy and yourself into this pit from which you don't get anything. 
							\item It's going to be reversed---you are in fact going to get much \textit{more} than you put in. 
							\item Not only will you come out \textit{even}, you will come out \textit{ahead}. You will get \textit{more}. 
							\item And that goes against our experience, it goes against this understanding we have of the world where it looks like you can't get more from something than you put into it. 
							\item But again, this is the promise.
						\end{compactitem}
					\item Alma too picks up on the idea at the end of his sermon on nourishing the seed so that it will grow into a tree, in \hyperref[ALMA-32]{Alma 32}.
						\begin{compactitem}
							\item He's about to finish this sermon, where he's got in mind as a background the parable of the sower that Christ says.
							\item Christ says the seeds fall in various places, and depending on the ground, they grow or don't grow.
							\item Alma focuses in on our role in preparing that ground and in taking care of the seed over time. 
							\item He's given us a lot of insights, and then he starts to wrap it up by saying, ``But if ye will nourish the word, yea, nourish the tree as it beginneth to grow by your faith with great diligence and with patience, looking forward to the fruit thereof, and it shall take root. And behold, it shall be a tree springing up unto everlasting life'' (\hyperref[ALMA-32-41]{Alma 32:41}).
							\item Notice that he uses the same language to describe the tree that Christ uses to describe the living water: it's so energetic, so active, that it's ``springing up,'' like a geyser or something.
							\item And then he goes on: ``And because of your diligence and your faith and your patience with the word in nourishing it that it may take root in you, behold, by and by ye shall pluck the fruit thereof, which is most precious, which is sweet above all that is sweet, and which is white above all that is white, yea, and pure above all that is pure; and ye shall feast upon this fruit even until ye are filled, that ye hunger not, neither shall ye thirst'' (\hyperref[ALMA-32-42]{Alma 32:42}).
							\item He makes the same point: it's not just desirable, but we can ``feast'' on it ``until we are filled.'' 
							\item As much as we want, it's ours. We put in our part---which, let's be honest, is going to be pretty meager---but then it just flows back to us endlessly.
							\item In that parable of the sower, Christ makes clear that what is produced in this process goes far beyond what we put in---he says at \hyperref[MATTHEW-13-23]{Matthew 13:23}, ``But he that received seed into the good ground is he that heareth the word, and understandeth it; which also beareth fruit, and bringeth forth, some an hundredfold, some sixty, some thirty.'
							\item The exact amount, like a hundred or sixty times more or whatever, is not as important as the fact that you get \textit{more} than what you put in. 
							\item There's no energy loss here, no waste heat, no thing like that.
							\item Just light and power and knowledge flowing back into our lives in return for what we give.
						\end{compactitem}
				\end{compactitem}
				
			\item Christ is the inexhaustible source of energy that pushes back against all those physical constraints that I talked about earlier---he is the thing that changes all the calculations, and that completely reverses this situation of energy loss in our lives.
				\begin{compactitem}
					\item He is the fountain that never runs dry; he has an unlimited supply of energy to give us, because he is the \textit{source} of that energy.
					\item DC 88
					\item The Atonement is this miraculous event where all these physical things and all these spiritual things collide.
					\item DC 29, says that those things are ultimately the same-ish
				\end{compactitem}
				
			\item Other ideas
				\begin{compactitem}
					\item Being led one step at a time, because I can't process anything farther ahead than that. I can't 
				\end{compactitem}
		
		\end{compactitem}
	
	% ----------------------
	\subsection{25 August 2024: Sacrament meeting talk on ``Words Matter''}
	% ----------------------
		\label{TALKS-25-08-2024}
		
		Ron Caldwell-Andrews asked Alicia and me to speak on this day in sacrament meeting. He said we could talk about whatever but suggested Rasband's April 2024 talk ``Words Matter'' as a base.
		
		\begin{compactitem}
		
			\item Hi everyone. My name is Bryce.
				\begin{compactitem}
					\item We are long-time members of this ward, which means we've been here for like five months, since the ward started existing.
					\item We have lived in Virginia for five years. We moved here from North Carolina.
					\item I'm actually disappointed I'm speaking after my wife, because I know from experience that people usually like her talks much more than they like mine.
					\item I teach at Southern Virginia University in BV, and my classes are pretty universally disliked, so I know the students are super glad to see me up here. If my wife taught the classes people would probably like them better.
					\item Now in those classes, I normally use music to signal to the students when it's appropriate for them to have certain emotional reactions.
					\item So when it's an exciting moment, I'll play an exciting song, just in the background so my students will know, ``Oh, I should listen now.'' Because why would you listen to all the other stuff if you don't know it's going to be exciting?
					\item Now in class when we come to a vulnerable moment, a moment where I'm sharing personal things, like pictures of me in high school or Blink 182 lyrics or something, then I always play the Goo Goo Dolls song ``Iris.'' You might remember.
					\item So, since I'm not going to be able to play that song on repeat as I share these remarks with you, I hope that you'll hear it in your heart.
					\item And if you don't remember that one, other Goo Goo Dolls songs would also be acceptable, like maybe ``Slide'' or ``Black Balloon'' or something.
				\end{compactitem}
				
			\item Anyway, today I'll be talking to you about the words of Christ. I'm going to go through a few verses to try to bring out some stuff that I've just been noticing recently.
			
			\item I want to start with a question: \textbf{Why is it possible for our testimony of Christ to grow and shrink over time?}
				\begin{compactitem}
					\item I've been thinking about this question in comparison to the knowledge or understanding of people from history, for example.
					\item Take Martin Luther King Jr. I know he existed, I know some facts about him, I've heard some of his speeches and read some of his writings.
					\item I've got knowledge of the guy---you could say I have a testimony of him, in that I have this belief in his existence and in the things he did.
					\item But notice that this knowledge, or testimony, doesn't ever change over time. It's not like my beliefs about King are sometimes stronger or sometimes weaker. I'm always the same level of confident that he existed, that he did what he did, and that it was important for the United States and for the world.
					\item But our belief in Christ is \textit{not} like this. At least, maybe it's different for you guys, but for me I do experience this waxing and waning of confidence, and this occasional growing and shrinking of my knowledge of Christ and the atonement.
					\item So why is it different? Why doesn't a belief in Christ function the same way as a belief in some important historical figure, like Martin Luther King Jr.?
				\end{compactitem}
				
			\item Well right off the bat we could come up with some reasons.
				\begin{compactitem}
					\item First, a belief in Christ is not just a list of facts that you say \textit{yes} to---like if I believe in Christ, it isn't just that I say, well, what it means for me to believe in Christ is, to know that he existed, to know that he did certain things, to know that he lived in certain ways.
					\item Those are facts, and they are real, but that isn't what it means to have a belief in Christ.
					\item Believing in Christ is actually a matter of having a \textit{relationship} with him.
					\item But I could have a relationship with MLK, in that I could be inspired by him, I could study his life, and he could motivate me to be a better person. He still has this effect on people.
					\item So having a belief in Christ has to be more than that---it has to be more than me being inspired by him or looking at him as a wise teacher or a great ethical mind or something like that.
					\item There must be something different or something additional about my relationship with him.
					\item And that difference must be able to \textit{make a difference} in \textit{me}, so that by having that relationship, I become able to do things that I could not do before.
					\item In particular, my relationship with Christ has to make the difference that exaltation is now possible for me, when it wasn't possible before.
					\item Because as good as King Jr. was, a belief in him isn't going to save anyone.
				\end{compactitem}
				
			\item One of the ways Christ makes this difference in our relationship is through his ``words,'' and the affect that his words have on us.
				\begin{compactitem}
					\item But again, it's not like the words of a historical figure, like MLK; for as inspiring and powerful as his words are, they are not the same as the word of Christ, or the word of God.
					\item King Jr. spoke words and people heard them, or they read them now.
					\item That is not really what we mean by the word of Christ, or the word of God. 
					\item It is a different kind of word that has a different kind of effect on us.
					\item And that's what I want to talk about today---to understand some of these differences with the word of God and the words of Christ, and to see why those differences make it so that our relationship with Christ necessarily has to be different from our other relationships.
					\item And this will help us answer the question from before, of why it's possible for our testimony of Christ to grow and shrink over time.
				\end{compactitem}
				
				
			\item Let's start by reading some scriptures about the \textit{word}: the \textit{word of God}, or the \textit{word of Christ}.
				\begin{compactitem}
							% Start with 1 Nephi 11:21--25, to establish that the word leads us to the fountain of living waters; then go to DC 93:19 to get the bit about how the word tells us \textit{how} to do something; this is the sense in which it leads us to the tree---we follow it by understanding Christ and by worshipping him, by making him a part of our lives.
							% Then go to John 4 to talk about the living waters, which the word leads us to. God's love, as the fountain of water and the tree, can then go to work on us.
							% In this sense we are nourished by the word, which is leading us to this love; both Jacob (\hyperref[JACOB-6-7]{Jacob 6:7}) and Moroni (\hyperref[MORONI-6-4]{Moroni 6:4}) speak of being ``nourished'' by the ``good word'' of God.
							% But there is another sense that the scriptures talk about the word of God, or about Christ's words---they talk about it like 
					\item \hyperref[1NEPHI-11-21]{1 Nephi 11:21--25}
						\begin{compactitem}
							\item Notice the specific order in which things happen in these verses. Why does it happen in this order? Why does the angel ask about the meaning of the tree just after Nephi has seen Mary holding the child Jesus?
							\item One answer is that Nephi's attention has just been drawn to Christ, and so he's ready to appreciate the fact that Christ is the most significant manifestation of God's love.
							\item Then Nephi sees a bunch of people worshipping Christ as he goes into the world. And then Nephi, without any prompting from the angel, sees and understands that the iron rod is the word of God, and that it leads to the ``fountain of living waters or to the tree of life,'' which we already know are supposed to represent God's love.
							\item So it's after seeing these people who are having a \textit{relationship} with Christ that Nephi tells us something about the word. In his thought the word is connected to the relationship.
						\end{compactitem}
					\item \hyperref[DC93-19]{DC 93:19}
						\begin{compactitem}
							\item The ``sayings,'' or the words of Christ, have a specific purpose: to teach us \textit{how} to do something---to understand and know \textit{how} to worship, or how to perform those actions and be in that relationship in the right way.
							\item The people Nephi saw were already worshipping Christ, so they had already received enough of Christ's words to understand their relationship with him, and on that basis they understood ane knew how to worship. And they did it.
						\end{compactitem}
					\item \hyperref[JOHN-4-14]{John 4:14}
						\begin{compactitem}
							\item How do you get this water?
							\item The word of God leads you to it, as Nephi taught; it leads you to \textit{do} something, and to get this active, living force inside of you so it can have an effect on \textit{who you are} and \textit{what you become}.
						\end{compactitem}
					\item So we've got some stuff on the table here about the word, and how it leads us to Christ, and more specifically it leads us to have a \textit{relationship} with Christ, so that we can get the waters inside of us.
						\begin{compactitem}
							\item There's a verse I didn't read in John that tells us something about that water.
							\item In \hyperref[JOHN-4-10]{John 4:10}, Christ describes the water as ``living'' water.
							\item That idea, that the water is \textit{living}, will help us pick up a few more nuances here that will take us back to the question we started with, about our testimony of Christ.
						\end{compactitem}
					\item \hyperref[HEBREWS-4-12]{Hebrews 4:12}
						\begin{compactitem}
							\item (Compare also \hyperref[HELAMAN-3-29]{Helaman 3:29} and \hyperref[DC6-2]{DC 6:2}.)
							\item Here we read that the word of God is \textit{quick}, which is kind of a strange thing to say.
							\item But ``quick'' is an old way of talking about life or being alive.
							\item And, in fact, the word translated as ``quick'' in Hebrews 4:12 is the same Greek word translated as ``living'' in John 4:10.
							\item The word itself is a \textit{living force}, and it also \textit{causes us to be more alive}.
							\item The ``living'' part is important. Why would it be ``living''? What does it mean that it's ``living''?
							\item Well, like we said, one of the meanings is that it causes us to be alive; by interacting with it, we are made more alive in some sense.
							\item But the fact that the word is living also means something else---it means that \textit{we} have to nourish \textit{it}. We have to nourish the word.
							\item We are nourished \textit{by it}, but we also have to nourish it ourselves in turn, or else it isn't going to grow; it won't become what it needs to be \textit{in us}.
							\item If we don't nourish it, nothing changes about the word itself, or about that living force.
							\item But we change, in that we give up the chance to have it be inside us and work inside us.
							\item And this bit about the word being alive helps us see something about the vision of the tree of life as well.
							\item It's not that the word as the iron rod leads you to a different location; the point is not to travel to a different place. The fact that it goes along a path is just a metaphor, or part of the vision.
							\item The word \textit{gets inside of us}, and as we nourish it through \textit{our actions}, and our \textit{relationship} with Christ, it \textit{grows into} the tree; it \textit{becomes} the fountain of living water, \textit{in us}. The tree is inside of us.
							\item Being able to go up to the tree of life and pick the fruit and eat it, to have the experience Lehi had, is a matter of getting the word inside of us and letting that tree grow in us, so that we become the kind of people for whom the atonement of Jesus Christ is the organizing principle of their lives.
							\item But because the word is a living thing, we have to take care of it; it's our job to make sure that this growth process happens in us. 
						\end{compactitem}
					\item \hyperref[ALMA-32-40]{Alma 32:40--43}
						\begin{compactitem}
							\item Let's finish off the scriptures by turning to Alma's commentary on the parable of the sower, where he ties together a bunch of these points at the end of Alma 32.
							\item Why do we have to ``nourish'' the word? Because it's alive! It's a living thing that has to be taken care of.
							\item And now we can finally go back to the question we started with.
							\item This is why it's possible for our testimony to grow and shrink over time: because that testimony is a \textbf{living thing}.
							\item And living things can grow and thrive, or they can shrink and waste away. Because life, or living things, are fighting against forces that want to disorganize them; it takes constant energy to keep a living thing alive, and even more energy for it to grow.
							\item This is why the ground in our hearts matters---we are trying to grow something. We are trying not just to have something stay alive, but for it to \textit{thrive}.
						\end{compactitem}
				\end{compactitem}
				
			\item I want to put the point a certain way but it's a bit unusual, so we need to explain a bit. I want to put it this way: Christ is as real to us as we want him to be.
				\begin{compactitem}
					\item Don't misunderstand what I mean here. The gospel isn't imaginary, it's not an illusion, it's not something we dream up or whatever.
					\item Christ is real, independently of us. Dana said this last week: he is a real person, and he remains real whether I believe in him or not.
					\item But we don't just need Christ to be real; that's true of lots of historical figures.
					\item We need him to be real \textit{for us}. And that is a very different matter.
					\item This is a matter of having a relationship with him, and of getting the word inside of us so it flowers and grows and thrives.
					\item \textbf{SEE ADDITIONAL MATERIAL BELOW IF I NEED IT.}
					\item All of this is real. Christ is real---the atonement is real. And in the end this is all that matters. 
				\end{compactitem}
				
			\item (Additional matreial if I need it)
				\begin{compactitem}
					\item This idea of ``Christ is as real to us as we want him to be'' is part of what Christ means in the Sermon on the Mount when he instructs us to ask, seek, and knock. You have to assume all of this is real \textit{to you} in order for those actions to make sense.
					\item Alma talks about this back in Alma 32, actually.
					\item Verse 34 is all about the effects that we will notice in our experience after we plant the word inside of us.
					\item Then we get \hyperref[ALMA-32-35]{verse 35}.
					\item This verse, and some of the other material here in Alma 32, is partly a response to Korihor. He accuses Alma of deceiving people.
					\item So Alma responds by talking about how no, we're not deceived people, and we ourselves aren't deluded either, because this is real. It is not a delusion. You are not deluded and you have not been deceived. It is all just as real as the fact that I am standing in front of you.
					\item You can tell it is real BECAUSE YOU CAN DISCERN IT!
					\item You have these beautiful moments with the spirit and the words of Christ, where it feels like the secrets of the universe are just opened to you and you feel like you're seeing things truly and clearly without any obstruction at all.
					\item But those moments are so rare, so achingly rare, and when the spirit recedes a bit and you have to come back down into the world it feels like immediately the scales of darkness just cover up your eyes and you're left to slog through this lonely waste again, and even when you come to church you might feel like those scales of darkness actually thicken because you might not feel totally welcome here, or you might feel like people don't understand you, or that your feelings and viewpoints on certain issues just don't fit with what it seems like the norms are.
					\item And so you feel like you're forced to make this choice: there are these moments when it feels like you're seeing everything so clearly, like the words of Christ are just flowing like a river into your mind, but there are many, many \textit{more} moments where it feels like everything is cloudy and gray and just you can't see anything because everything is such a mess.
					\item And it feels like you're forced to make this choice---you're forced to decide what is real. And if you haven't had any of those moments of clarity recently, then it's even more difficult to make this choice because the clarity is lost to so many years in the past.
					\item But I am here to tell you that those moments when you feel enlightened, and you feel the words of Christ inside you as a well of water springing up, \textit{those} are the moments when you are seeing most clearly; \textit{those} are the moments when you are grasping the truth.
					\item \textit{Those} are the moments when the true nature of things is revealed to you most clearly. That is the vision you have to try to hold onto.
					\item And one of the greatest challenges of mortality, at least for me, is to fix in my memory and in my vision that clarity, when everything else around has gone gray.
					\item Sometimes, unfortunately, we are at a stage of our spiritual lives when the fruit of the tree just seems to be out of season; it's a time when we have to wait for the fruit to appear.
					\item Alma is really specific about this at the end of Alma 32---he says there are times when you will have to wait for the tree to bring forth fruit.
					\item This is really frustrating sometimes, and even sad, because you're desperate to eat that fruit again.
					\item But the atonement is real, and the atonement is the best word you can plant inside you; it will never fail to bring forth the fruit we want. Sometimes we just have to wait. The right thing to do is be patient.
					% \item I don't know why the institutional church is the way it is; frankly the less I know about it the better. But that does not change anything about the reality of the atonement.
				\end{compactitem}
						
		\end{compactitem}
				
			% Christ---what goes in and out of us; Matthew 15:10--20
				% It's about \textit{actions}.
				
			% Christ's words to us: 2 Nephi 32:1--6
				% Christ is speaking \textit{to us}. 
				% This is part of what make the scriptures so interesting---it's like a personalized set of messages for each person. Even though the words printed on the page do not change, the meaning does for us, as our lives and our circumstances and our knowledge changes.
			
			% We tend to think about whatever we put into our heads most frequently. If we're filling our heads with the words of Christ, we're going to think about them, and they're going to be more likely to have an impact on our action.
				% And the key is action. It's action that has the greatest affect on us, in terms of change---habits, either creating them, breaking them, or maintaining them, change us.
				% We become what we think about, but only by means of action. 
				
			% There is a reason that a remission of sin is connected to obeying the commandments, or in other words, to living righteously. It's because actions demonstrate that we are allowing ourselves to work with Christ, and allowing the Holy Ghost to purify us. Just uttering the words ``Forgive me'' is just words.
				% What about Alma and other conversions in the BOM, though?
				% Making the effort is enough. 
				
			% I studied the brain when I was in graduate school, and the brain is really interesting, people love to learn about it. But neuroscientists are human beings, so they like to kind of hype up the things they learn about.
				% Here's one of those hyped-up claims. Sometimes people will say that doing a particular thing ``changes your brain.'' Like they'll say, learning a language changes your brain, or learning how to play a sport changes your brain.
				% Those claims aren't false, but they also aren't as important or as interesting as they seem, because in reality, \textit{everything} you do, all the time, changes your brain.
				% Right now, you listening to me, that's changing your brain. If it weren't then you wouldn't be able to understand me, or in the case of my undergraduates, you wouldn't be able to ignore me. Your brain \textit{has} to change for you to experience \textit{anything}.
				% Sure, exercise changes my brain. But do you know what else does? Eating this entire box of donuts. So if changing your brain is what you're going for, you've actually got a couple of options here.
				% What we are about are long-term changes that are resilient.
	
	% ----------------------
	\subsection{28 April 2024: EQ discussion on the talk ``Covenant Belonging''}
	% ----------------------
		\label{TALKS-28-04-2024}
		
		The talk for this discussion is Gong's October 2019 talk ``Covenant Belonging,'' available \href{https://www.churchofjesuschrist.org/study/general-conference/2019/10/41gong?lang=eng}{here}.
		
		\begin{compactitem}
		
			\item Introducing Gong's talk ``Covenant Belonging''
				\begin{compactitem}
					\item This is the talk that bishop has asked us to read.
					\item If you haven't read it, it's pretty interesting, and it's worth taking a look.
					\item Today we won't be looking at all or even most of this---we're going to focus on a few remarks that Gong makes in a few different places in order to push us into a larger discussion about Christ's role in the plan of salvation.
					\item So let's take a look at a few passages from the talk.
						\begin{compactitem}
							\item First, we'll look at the seventh paragraph, which starts, ``God's ordinances and covenants...'' Could someone read that paragraph for us?
							\item What really sticks out to me here is the last sentence, about redemption and salvation; they are possible because of Christ's atonement.
							\item Then, two paragraphs later, Gong expands on that---he says, ``Covenant belonging centers in Jesus Christ as `mediator of the new covenant'.''
							\item Christ is able to be this mediator because he performed the atonement. Again, the key here is his sacrifice; the gospel is nothing without it. Everything depends on that.
							\item Let's read one more bit, a part about half way through the talk, in a paragraph that begins, ``The Book of Mormon invites each of us...''
							\item Let's focus on the last sentence.
							\item Gong says that ``we can come by covenant to belong with God,'' and by the idea of ``covenant belonging,'' he seems to mean the fact that we are join a community of followers of Christ when we make covenants. 
							\item I want to explore this idea of making covenants in more detail, and I want to do it in order to approach a different question, that I think is fundamental for our understanding of the gospel.
						\end{compactitem}
				\end{compactitem}
		
			\item That question is, what is the relationship between the atonement and the priesthood?
				\begin{compactitem}
					\item We are not going to be able to answer this fully today; I am confident that I'm not able to answer it fully yet.
					\item I hope one day to be able to.
					\item But we can make some headway on it by looking at some scriptures.
				\end{compactitem}
			
			\item Let's begin by reading some scriptures about people making covenants.
				\begin{compactitem}
					\item Gong is talking about covenant belonging, and he seems to be referring to our making covenants with God and with each other. 
					\item We can also think of covenant belonging in another sense, one that appears in different parts of the scriptures. So let's take a look at some.
					\item \hyperref[MOSIAH-18-20]{Mosiah 18:20--22}
						\begin{compactitem}
							\item sdf
						\end{compactitem}
					\item \hyperref[3NEPHI-9-17]{3 Nephi 9:17}
						\begin{compactitem}
							\item sdf
						\end{compactitem}
					\item \hyperref[MORONI-7-48]{Moroni 7:48}
						\begin{compactitem}
							\item sdf
						\end{compactitem}
					\item \hyperref[DC11-30]{DC 11:30}
						\begin{compactitem}
							\item sdf
						\end{compactitem}
					\item \hyperref[DC25-1]{DC 25:1}
						\begin{compactitem}
							\item This one mentions women too; it isn't just for men.
						\end{compactitem}
					\item \hyperref[DC34-3]{DC 34:3}
						\begin{compactitem}
							\item sdf
						\end{compactitem}
					\item \hyperref[DC35-2]{DC 35:2}
						\begin{compactitem}
							\item Important because it mentions the in-dwelling relationship.
							\item See also \hyperref[DC50-43]{DC 50:43}.
						\end{compactitem}
					\item \hyperref[DC39-4]{DC 39:4}, \hyperref[DC42-52]{DC 42:52}, \hyperref[DC45-8]{DC 45:8}, etc.
					\item The key here is that certain people ``\textit{become} the sons (or children) of God.''
						\begin{compactitem}
							\item But why would that be?
							\item Aren't we alread the children of God? Why do the scriptures talk about us becoming the children of God, if we already are?
							\item This language is referring to a new covenant relationship with God.
						\end{compactitem}
					\item But King Benjamin, after he gives his sermon, is more specific than this, and teaches us some other very important things.
						\begin{compactitem}
							\item Let's begin reading at \hyperref[MOSIAH-5-6]{Mosiah 5:6}.
							\item (Two important points: first, we become the children of Christ because of our covenant relationship \textit{with him}; second, we are spiritually reborn \textit{through him}; in other words, it is by virtue of \textit{his atonement} that we are able to be spiritually reborn in this transformation.)
							\item (Read 4:2--3 and 5:1--8 to expand on these points.)
							\item (Then read Moses 6:66--68.)
							\item (Then get into the connection between the covenants we make and entering into the order of the priesthood.)
							\item (Then get into DC 107.)
							\item (Then get into, why is it called that? Why is the priesthood after the order of the son of God? Well, it is that way for us; we access the power of our heavenly parents through Christ; he is the mediator of the new covenant. Get into mediator as \gr{μεσίτης}, ``mediator, arbitrator, guarantor.'')
							\item Christ is the channel through which we get access to the priesthood power that makes salvation possible.
							\item But what makes it possible for Christ to do this?
							\item The answer is, the atonement.
							\item Matthew 28; DC 93:16--17
							\item DC 88, Alma 13
						\end{compactitem}
				\end{compactitem}
			
			\item By carrying out the atonement, Christ becomes the central figure in the plan of salvation, as he was always intended to be.
			
			\item Scriptures
				\begin{compactitem}
					\item \hyperref[MATTHEW-28-18]{Matthew 28:18}, ``And Jesus came and spake unto them, saying, All power is given unto me	in heaven and in earth.''
						\begin{compactitem}
							\item Note the difference between \gr{ἐξουσία} and \gr{δύναμις}.
							\item Note the alternate translations ``has been given'' or ``was given''; ``has been given'' is far and away the most common translation in other versions.
							\item NICNT has ``All authority in heaven and on earth has been given to me.''
							\item NIGTC has ``All authority in heaven and on earth has been given to me.''
							\item This goes well with JS talking about salvation as triumphing over all enemies, the last of which is death; see quotes from JS below for 14 and 21 May 1843.
						\end{compactitem}
					\item \hyperref[HEBREWS-12-22]{Hebrews 12:22--24}, ``But ye are come unto mount Sion, and unto the city of the living God, the heavenly Jerusalem, and to an innumerable company of angels, To the general assembly and church of the firstborn, which are written in heaven, and to God the Judge of all, and to the spirits of just men made perfect, And to Jesus the mediator of the new covenant, and to the blood of sprinkling, that speaketh better things than that of Abel.''
					\item \hyperref[ALMA-13-5]{Alma 13:5}, ``or in fine, in the first place they were on the same standing with their brethren--thus this holy calling being prepared from the foundation of the world for such as would not harden their hearts, being in and through the atonement of the Only Begotten Son which was prepared---''
					\item \hyperref[MORONI-10-33]{Moroni 10:33}, ``And again, if ye by the grace of God are perfect in Christ and deny not his power, then are ye sanctified in Christ by the grace of God through the shedding of the blood of Christ, which is in the covenant of the Father, unto the remission of your sins, that ye become holy, without spot.''
					\item \hyperref[DC76-56]{DC 76:56}, ``They are they who are priests and kings, who have received of his fulness, and of his glory;''
						\begin{compactitem}
							\item Fulness of what? Of the Father.
							\item ``And again, we saw the terrestrial world, and behold and lo, these are they who are of the terrestrial, whose glory differs from that of the church of the Firstborn who have received the fulness of the Father, even as that of the moon differs from the sun in the firmament'' (\hyperref[DC76-71]{DC 76:71})
							\item ``These are they who receive of the presence of the Son, but not of the fulness of the Father'' (\hyperref[DC76-77]{DC 76:77})
							\item ``Ye who are \hyperref[QUICKEN]{quickened} by a portion of the celestial glory shall then receive of the same, even a fulness'' (\hyperref[DC88-29]{DC 88:29})
							\item See also \hyperref[DC93-19]{DC 93:19}.
						\end{compactitem}
					\item \hyperref[DC76-57]{DC 76:57}, ``And are priests of the Most High, after the order of Melchizedek, which was after the order of Enoch, which was after the order of the Only Begotten Son.''
					\item \hyperref[DC84-19]{DC 84:19--20}, ``And this greater priesthood administereth the gospel and holdeth the key of the mysteries of the kingdom, even the key of the knowledge of God. 20 Therefore, in the ordinances thereof, the power of godliness is manifest.''
					\item \hyperref[DC88-5]{DC 88:5--13}, ``Which glory is that of the church of the Firstborn, even of God, the holiest of all, through Jesus Christ his Son--6 He that ascended up on high, as also he descended below all things, in that he comprehended all things, that he might be in all and through all things, the light of truth; 7 Which truth shineth. This is the light of Christ. As also he is in the sun, and the light of the sun, and the power thereof by which it was made. 8 As also he is in the moon, and is the light of the moon, and the power thereof by which it was made; 9 As also the light of the stars, and the power thereof by which they were made; 10 And the earth also, and the power thereof, even the earth upon which you stand. 11 And the light which shineth, which giveth you light, is through him who enlighteneth your eyes, which is the same light that quickeneth your understandings; 12 Which light proceedeth forth from the presence of God to fill the immensity of space--13 The light which is in all things, which giveth life to all things, which is the law by which all things are governed, even the power of God who sitteth upon his throne, who is in the bosom of eternity, who is in the midst of all things.'' 
						\begin{compactitem}
							\item More connections here between Christ experiencing all things, going below all things, and thus being connected to all things through the priesthood.
						\end{compactitem}
					\item \hyperref[DC88-19]{DC 88:19--22}, ``For after it hath filled the measure of its creation, it shall be crowned with glory, even with the presence of God the Father; 20 That bodies who are of the celestial kingdom may possess it forever and ever; for, for this intent was it made and created, and for this intent are they sanctified. 21 And they who are not sanctified through the law which I have given unto you, even the law of Christ, must inherit another kingdom, even that of a terrestrial kingdom, or that of a telestial kingdom. 22 For he who is not able to abide the law of a celestial kingdom cannot abide a celestial glory.''
					\item \hyperref[DC93-16]{DC 93:16--17}, ``And I, John, bear record that he received a fulness of the glory of the Father; And he received all power, both in heaven and on earth, and the glory of the Father was with him, for he dwelt in him.''
						\begin{compactitem}
							\item It is only \textit{after} the atonement that Christ receives this power.
						\end{compactitem}
					\item \hyperref[DC107-1]{DC 107:1--6}, ``THERE are, in the church, two priesthoods, namely, the Melchizedek and Aaronic, including the Levitical Priesthood. 2 Why the first is called the Melchizedek Priesthood is because Melchizedek was such a great high priest. 3 Before his day it was called the Holy Priesthood, after the Order of the Son of God. 4 But out of respect or reverence to the name of the Supreme Being, to avoid the too frequent repetition of his name, they, the church, in ancient days, called that priesthood after Melchizedek, or the Melchizedek Priesthood. 5 All other authorities or offices in the church are appendages to this priesthood. 6 But there are two divisions or grand heads---one is the Melchizedek Priesthood, and the other is the Aaronic or Levitical Priesthood.''
					\item \hyperref[MOSES-6-62]{Moses 6:62--68}, ``And now, behold, I say unto you: This is the plan of salvation unto all men, through the blood of mine Only Begotten, who shall come in the meridian of time. 63 And behold, all things have their likeness, and all things are created and made to bear record of me, both things which are temporal, and things which are spiritual; things which are in the heavens above, and things which are on the earth, and things which are in the earth, and things which are under the earth, both above and beneath: all things bear record of me. 64 And it came to pass, when the Lord had spoken with Adam, our father, that Adam cried unto the Lord, and he was caught away by the Spirit of the Lord, and was carried down into the water, and was laid under the water, and was brought forth out of the water. 65 And thus he was baptized, and the Spirit of God descended upon him, and thus he was born of the Spirit, and became quickened in the inner man. 66 And he heard a voice out of heaven, saying: Thou art baptized with fire, and with the Holy Ghost. This is the record of the Father, and the Son, from henceforth and forever; 67 And thou art after the order of him who was without beginning of days or end of years, from all eternity to all eternity. 68 Behold, thou art one in me, a son of God; and thus may all become my sons. Amen.''
					\item \hyperref[JST-OTR1-GENESIS-14]{Joseph Smith Translation Genesis 14}, ``And Melchizedeck lifted up his voice and blessed Abram. Now Melchizedeck was a man of faith who wrought righteousness and when a child he feared God and stoped the mouths of Lions and quenched the violence of fire and thus having been approved of God he was ordained \sout{<a>} a high Preist after the \hyperref[ORDER]{order} of the covenent which God made with Enoch it being after the \hyperref[ORDER]{order} of the Son of God which \hyperref[ORDER]{order} came not by man nor the will of men \sout{but of} neither by father nor Mother neither by begining of days nor end of years but of God and was delivered unto men by the calling of his own voice according to his own will unto as many as beleived on his name for God having sworn unto Enoch and unto his seed with an oath by himself that every one being ordained after this \hyperref[ORDER]{order} and calling should have power by faith to break Mountains to divide the seas to dry up watters to turn them out of their course to put at defience the armies of nations to divide the earth to break every band to stand in the preasence of God to do all things according to his will according to his command subdue principalities and powers and this by the will of the Son of God which was from before the foundation of the world and men having this faith coming up unto this \hyperref[ORDER]{order} of <God> were translated and taken up into heaven. And now Melchizedeck was a Priest of this \hyperref[ORDER]{order} therefore he obtained peace in Salam and was called the prince of peace and his people wraught righteousness and obtained \sout{heaven} heaven and sought for the City of Enoch which God had before taken seperating it from the earth having reserved it unto the latter days or the end of the world and hath said and sworn with an oath that the heavens and the earth should come together and the sons of God should be tryed so as by fire and this Melchizedeck having thus established righteousness was called the King of heaven by his people or in other words the King of peace and he lifted up his voice and he blessed Abram being the high Priest and the keeper of the store house of God him whom God had appointed to receive tithes for the poor wherefore Abram paid unto him tithes of all that he had of all the riches which he posessed which God had given him more \sout{that} than that which he had need and it came to pass that God blessed Abram and gave unto him riches and honour and lands for an everlasting possession according to the covenent which he had made and according to the blessing wherewith Melchizedeck had blessed him.''
				\end{compactitem}
			
			\item JS and BY quotations
				\begin{compactitem}
					\item JS: ``There are two priesthoods spoken of in the scriptures, viz, the Melchisadeck and the Aaronic or Levitical Altho there are two priesthoods, yet the Melchisadeck priesthood comprehends the Aaronic or Levitical priesthood and is the Grand head and holds the hig[h]est Authority which pertains to the priesthood--the keys of the Kingdom of God in all ages of the world to the latest posterity on the earth--and is the channel through which all knowledge, doctrine, the plan of salvation and every important \sout{matter} truth is revealed from heaven. Its institution was prior to ``the foundation of this earth or the morning stars sang together or the sons of God shouted for joy,'' \sout{and} it is the highest and holiest priesthood and is after the order of the Son [of] God, and all other priesthoods are only parts, ramifications, powers and blessings belonging to the same and are held controlled and directed by it. It is the Channel through which the Almighty commenced revealing his glory at the beginning of the creation of this earth and through which he has continued to reveal himself to the children of men \sout{and} to the present time and through which he will make known his purposes to the end of time---''
						\begin{compactitem}
							\item \hyperlink{EMS-1840-JS-CIR-5-OCT-1840-b}{5 October 1840}
						\end{compactitem}
					\item JS: ``The offering of sacrifice \sout{is} has ever been connected, and forms a part of the <duties of the> priesthood. It began \sout{which} with the priesthood and will be continued untill after the coming of christ from generation to generation''
						\begin{compactitem}
							\item \hyperlink{EMS-1840-JS-CIR-5-OCT-1840-w}{5 October 1840}
						\end{compactitem}
					\item JS: ``All priesthood is Melchizedeck; but there are different portions or degrees of it.''
						\begin{compactitem}
							\item \hyperlink{EMS-1841-JS-5-JAN-1841-WC-b}{5 January 1841}
						\end{compactitem}
					\item JS: ``Salvation is nothing more or less than to triumph over all our enemies \& put them under our feet \& when we have power to put all enemies under our feet in this world \& a knowledge to triumph over all evil spirits in the world to come then we are saved as in the case of Jesus he was to reign untill he had put all enemies under his feet \& the last enemy was death''
						\begin{compactitem}
							\item \hyperlink{EMS-1843-JS-14-MAY-1843-A-WW-a}{14 May 1843}
						\end{compactitem}
					\item JS: ``\uline{salvation}. is for a man to be seveed from all his enemies.--until a man can triumph over death. he is not saved. knowlidge will do this''
						\begin{compactitem}
							\item \hyperref[EMS-1843-JS-21-MAY-1843-WR]{21 May 1843}
							\item Knowledge of what? Of Christ: ``I give unto you these sayings that you may understand and know how to worship, and know what you worship, that you may come unto the Father in my name, and in due time receive of his fulness'' (\hyperref[DC93-19]{DC 93:19}).
						\end{compactitem}
					\item JS: ``What is Salvation. Salvation is for a man to be Saved from all his enemies even our last enemy which is death''
						\begin{compactitem}
							\item \hyperref[EMS-1843-JS-21-MAY-1843-HC]{21 May 1843}
						\end{compactitem}
				\end{compactitem}
		
		\end{compactitem}
	
	% ----------------------
	\subsection{10 March 2024: EQ discussion on the talk ``The Prodigal and the Road That Leads Home''}
	% ----------------------
		\label{TALKS-03-10-2024}
		
		The talk for this discussion is Uchtdorf's October 2023 ``The Prodigal and the Road That Leads Home,'' available \href{https://www.churchofjesuschrist.org/study/general-conference/2023/10/45uchtdorf?lang=eng}{here}.
		
		\begin{compactitem}
		
			\item Introduce Uchtdorf's talk about the parable he's going to speak about
			
			\item Let's talk about our traditional way of referring to this parable. What does the word ``prodigal'' mean?
				\begin{compactitem}
					\item 1828 dictionary for the adjective ``prodigal'':
						\begin{compactitem}
							\item[1] Given to extravagant expenditures; expending money or other things without necessity; profuse, lavish; wasteful; not frugal or economical.
							\item[2] Profuse, lavish; expended to excess or without necessity.
							\item[3] Very liberal; profuse.
						\end{compactitem}
					\item This parable complements and extends the ones that are given before it, all circling back to certain attitudes held by the Pharisees in the opening verses of the chapter.
				\end{compactitem}
				
			\item There are three characters in this parable, and we can divide it into two rough parts.
				\begin{compactitem}
					\item The character that appears in both parts is the father. The emphasis is on him.
					\item We, however, are not the father; we are one or both of the sons, at different times.
				\end{compactitem}
				
			\item On the son ``coming to himself,'' or coming to his senses.
				\begin{compactitem}
					\item How have you experienced this ``coming to yourself'' or ``coming to your senses'' in the past?
					\item Were there particular events that caused you to ``come to yourself''?
						\begin{compactitem}
							\item In the story, the younger son is with the ``swine,'' or pigs.
							\item These were an abomination to the Jewish audience for this parable (\hyperref[LEVITICUS-11-7]{Leviticus 11:7}).
						\end{compactitem}
					\item We all depart from the path at different times and for different reasons. \textit{Is there any way of accelerating this process of ``coming to yourself'' when you depart?}
				\end{compactitem}
		
			\item In what ways do we manifest the attitude of the older brother to our family, friends, or other people around us?
			
			\item This parable has some important teachings for us about the difference between merit and grace.
				\begin{compactitem}
					\item \textbf{In what ways do we demand that God deal with us according to what we deserve?}
					\item Compare the use of ``merit'' or ``merits'' in the restoration scriptures.
						\begin{compactitem}
							\item Merits of Christ: \hyperref[2NEPHI-2-8]{2 Nephi 2:8}, \hyperref[2NEPHI-31-19]{2 Nephi 31:19}, \hyperref[ALMA-24-10]{Alma 24:10}, \hyperref[HELAMAN-14-13]{Helaman 14:13}, \hyperref[MORONI-6-4]{Moroni 6:4}, \hyperref[DC3-20]{DC 3:20}
							\item Our merits, or lack thereof: \hyperref[MOSIAH-2-19]{Mosiah 2:19}, \hyperref[ALMA-22-14]{Alma 22:14}
						\end{compactitem}
					\item Repentance is a change of mind, a forward-looking return to the truth in order to move your life in a different direction.
						\begin{compactitem}
							\item It is not, and can't be, about balancing the scales.
							\item We earn very little by obeying the commandments (Mosiah 2:24 below).
						\end{compactitem}
					\item Some scriptures
						\begin{compactitem}
							\item \hyperref[2NEPHI-10]{2 Nephi 10}: lay aside sin; reconcile yourselves to the will of God
							\item \hyperref[2NEPHI-25-23]{2 Nephi 25:23}: be reconciled to God; it is by grace we are saved, after all we can do
							\item \hyperref[MOSIAH-2-24]{Mosiah 2:24}: God paying us for obeying the commandments, but we haven't earned much when we do that.
						\end{compactitem}
				\end{compactitem}
				
			\item The irony is that we can be truly ``prodigal,'' or experience bounty without end, only through Christ.
				\begin{compactitem}
					\item \hyperref[2NEPHI-26-25]{2 Nephi 26:25}: ``Behold, doth he cry unto any, saying: Depart from me! Behold, I say unto you: Nay. But he saith: Come unto me, all ye ends of the earth; buy milk and honey without money and without price.''
				\end{compactitem}
		
		\end{compactitem}
	
	% ----------------------
	\subsection{11 February 2024: EQ discussion on the talk ``Praise the Man''}
	% ----------------------
		\label{TALKS-02-11-2024}
		
		This discussion is roughly based around Russell Ballard's talk ``Praise to the Man,'' available \href{https://www.churchofjesuschrist.org/study/general-conference/2023/10/41ballard?lang=eng}{here}.
		
		\begin{compactitem}
		
			\item Ballard gave this talk in October 2023, at general conference, which was very soon before his death.
				\begin{compactitem}
					\item If you remember the talk or heard it recently, Ballard seems to apprehend pretty clearly that his time on earth was coming to an end, and in fact he did pass away just a few weeks later.
					\item I've always found these sorts of talks interesting, that people give right before they pass.
					\item In this case Ballard chose to speak on the topic of Joseph Smith as the prophet of the restoration, and he called his remarks ``Praise to the Man,'' which of course refers to the hymn.
					\item Having very recently read 2 Nephi 1--3, which is Lehi's final message, I invite you to think about what you would say in a final address, if you were to give one.
					\item Let's take a look at some of the things Ballard says.
				\end{compactitem}
				
			\item The Tyson Zone
				\begin{compactitem}
					\item Let's set aside Joseph Smith for a bit and we'll come back to it soon.
					\item I want to introduce you to an idea called the ``Tyson Zone.''
					\item (Introduce the Tyson Zone.)
					\item Celebrities enter the Tyson Zone.
				\end{compactitem}
				
			\item Joseph Smith and the Tyson Zone
				\begin{compactitem}
					\item I sometimes wonder if, for me, Joseph Smith has entered the Tyson Zone.
					\item It may be that there's nothing so bizarre that you could tell me about Joseph Smith that I wouldn't believe, at least initially.
					\item And when I start to think along these lines, I wonder to myself, what would I have to find out about Smith in order to \textit{not} believe he was a prophet?
					\item Is there anything I could find out that could convince me to change the view I have now? There must be, though I can't say for sure right now what those things are.
					\item In any case, not everyone feels the same way as Ballard about Joseph Smith, obviously.
					\item Now I want to ask all of you about this, but from a different angle. That angle comes from something Christ taught.
				\end{compactitem}
				
			\item In Matthew 7, Christ teaches us about false prophets, and gives us a test that we can use to identify false prophets.
				\begin{compactitem}
					\item I'd like to read this test and ask all of you what you think about it and how to apply it.
					\item So let's read from \hyperref[MATTHEW-7-15]{Matthew 7:15--20}.
					\item Open up to those verses---could someone read them for us?
					\item 
					\item The principle here is, by the fruits of alleged prophets, we will be able to identify them as false or true.
					\item So to tie back to some of the stuff I was talking about earlier: \hl{On the basis of what fruits do we make this judgment?}
						\begin{compactitem}
							\item A person's life produces many fruits. Joseph Smith's did too, just like all of ours---our moment-to-moment behaviors, but also our professional lives, our family lives, our hobbies, and many other things.
							\item \hl{So what fruits do we make this judgment on, and why those and not others?}
						\end{compactitem}
				\end{compactitem}
				
			\item The larger question is, \hl{what aspects of a prophet’s behavior are relevant to judging whether that person is truly a prophet or not?}
				\begin{compactitem}
					\item How should these criteria affect the way \textit{we} live our lives? 
						\begin{compactitem}
							\item \hyperref[REVELATION-19-10]{Revelation 19:10}, ``the testimony of Jesus is the spirit of prophecy.''
							\item JS uses this verse a few different times to talk about how, in a certain sense, \textit{everyone} is a prophet.
							\item For example, see \hyperlink{EMS-1838-JS-8-MAY-1838-b}{these questions and answers} from May 1838.
							\item See also \hyperlink{EMS-1839-JS-CIR-26-JUN-CIR-4-AUG-1839-A-i}{these remarks} from summer 1839.
							\item (\hyperlink{EMS-1843-JS-11-JUN-1843-WW-l}{These remarks} are about there always being prophets; not the best one to share in this discussion.)
							\item Joseph Smith commented on some of this in 1843.
							\item For example, in \hyperlink{EMS-1843-JS-21-MAY-1843-WR-b}{May} he talked about how he wasn't that much better than other people.
							\item And in connection with this, he once made a very interesting remark about righteousness. According to \hyperlink{EMS-1843-JS-21-MAY-1843-HC-b}{another account} of the same remarks in May, he said, ``Righteousness is not that which men esteem holiness. That which the world call righteousness I have not any regard for. To be righteous is to be just and merciful.''
							\item What does this mean for us?
							\item As we evaluate ourselves as ``prophets,'' or as those with the testimony of Jesus, in what way do we use either (1) holiness, or (2) worldly conceptions of righteousness, to stand in for true righteousness?
						\end{compactitem}
				\end{compactitem}
		
		\end{compactitem}
	
	% ----------------------
	\subsection{6 October 2023: Why I Believe devotional at the SVU institute}
	% ----------------------
		\label{TALKS-06-10-2023}
		
		(I took extensive notes in preparation for this talk. The actual version, which I unfortunately forgot to record, ended up being a mixture of a bunch of the the things I had written, with some new things. See my discussion of the experience \hyperref[SE-WHY]{here}.)
		
		\begin{compactitem}
		
			\item I'm grateful to the institute council for inviting me to speak to you today, though I'm self-conscious about being here.
				\begin{compactitem}
					\item It isn't because I don't want to speak to you---I do---it's because I myself am not a very good devotional \textit{attender}.
					\item I have a hard time paying attention to these sorts of talks.
					\item When it comes to these things I have the attention span of a little kid. I get distracted by everything.
					\item So if any of you out there are distracted right now, I know you're not hearing this because you're distracted, but I feel you, because I am you.
				\end{compactitem}
			
			\item I've been thinking for almost two weeks now about what I would say, and I've actually struggled to decide.
				\begin{compactitem}
					\item I've written out lots of stuff a few different times but none of it felt right.
					\item So earlier this week I asked some people in one of my classes and some other students too, about what things you'd be interested in hearing me talk about.
					\item They shared some really useful stuff with me and so I'm just going to try to aim for the things they mentioned.
					\item The truth is that I feel a lot of pressure because I want to give you something meaningful, and I'm sometimes bad at figuring out what the audience is interested in compared to what I'm interested in.
					\item I'm not a great speaker. This is not false modesty, I am actually bad at organizing talks and speeches.
					\item I would much rather write things. I started writing essays about the gospel recently just to try to sort my own ideas out and it's kind of a laborious process because writing takes a while.
					\item So fair warning, this may be a disorganized mess. I am particularly bad when the topic is so open-ended. That's really hard for me.
				\end{compactitem}
				
			\item The question is why I believe, and it actually feels a bit like the question is why I \textit{still} believe.
				\begin{compactitem}
					\item One of the students I spoke to actually said this---she said, you should talk about why you still believe, after studying these things and having this background.
					\item So let me answer that question as simply as possible. There are two reasons why I believe, or why I still believe.
					\item I believe because I've had many experiences with the Holy Ghost, and I don't see any reason to think those experiences aren't real. 
					\item I still believe because the restored gospel of Jesus Christ has given me clear, insightful answers to questions that I couldn't have received elsewhere.
				\end{compactitem}
				
			\item But let me tell you about where my belief started.
				\begin{compactitem}
					\item When I left on my mission in 2005, I had never received an answer in prayer about whether the Book of Mormon was true or whether the church was true.
					\item This was kind of upsetting to me because it felt like I really needed it to happen before I left.
					\item I had prayed about those things many, many times, but I never felt like I got an answer about any of it, and to this day I never have.
					\item Before my mission I had been praying a lot about this and I remember having this thought that was like, wait, why do I need an answer in prayer? I already believe the Book of Mormon is true. Why do I need an additional answer if I already believe?
					\item And this got rid of my worry about it. I think my view on the Book of Mormon has changed, but I've never worried about it.
					\item So I look back now and wonder, why did I serve? 
					\item My parents hadn't served, and they didn't pressure me to do it. In fact we almost never spoke about it. I just went because that's what you did when you graduated high school in Utah.
					\item (More background to me.)
					\item But I get out there in the mission, and I start having all these very powerful experiences with the Spirit.
					\item Not just feelings during lessons or whatever---on my occasions I felt as though a force outside myself was putting into my mind the words I was supposed to say.
					\item This happened mostly when I was speaking to other missionaries.
					\item The voice would enter my mind, and I would speak whatever words it gave me. Then the voice would stop, and I would just wait. However the other person was responding, I would just wait until the voice came back, and when it did, I would say the words it gave me, and then I would stop.
					\item Those experiences, probably more than anything else, convinced me that the Holy Ghost was a real thing that could communicate with me and give me knowledge and insight that I couldn't get any other way.
					\item (How this convinced me of Christ and the atonement.)
					\item So I came home from my mission with a strong belief based in many experiences with the Spirit.
				\end{compactitem}
				
			\item This belief got challenged pretty fast after my mission, though, because soon after I got home I started to get pretty severe anxiety.
				\begin{compactitem}
					\item (How I had learned on my mission to associate the Spirit with emotions and physical sensations.)
					\item (Anxiety mimicked all those physical sensations, so I could no longer tell the difference.)
					\item (Extremely difficult for me now to tell when my feelings are spiritually significant or not.)
					\item (Story of Alicia and of going to Duke, the two most important decisions of my adult life, I never received an answer in prayer to either of those things, and I prayed a LOT.)
					\item (Studying it out in our own minds.)
					\item (The result has been that I have disconnected my emotional life from my spiritual life, in a way.)
					\item This has been difficult sometimes, since I want to feel God's love, not just know about it intellectually. But that is hard for me.
				\end{compactitem}
				
			\item Since I have trouble taking my emotions as reflective of the Spirit, the thing that has replaced them in my spiritual life has been the scriptures, and the words of Joseph Smith and Brigham Young.
				\begin{compactitem}
					\item (Why I love the scriptures.)
					\item You learn more about the atonement by reading the Book of Mormon than by doing anything else.
					\item (Why I love JS.)
					\item (Why I love BY.)
				\end{compactitem}
				
			\item But in more recent times I've had a very different kind of question about my experiences with the Holy Ghost, including those ones on my mission where it felt like there were words being given to me.
				\begin{compactitem}
					\item And this question has to do with the brain.
					\item (My background in neuro---start with BYU, then Dennett, then Duke.)
					\item Learning about the brain is really fascinating. I love it and so do many other people. Everyone loves to hear about the brain. It instantly interests people.
					\item And it's like this universal thing you can study---you can study the brain and its relation to poetry, or music, or logical reasoning, or anything else you want. 
					\item And whatever you study with it, you seem to get these explanations for why people do what they do, or why they believe what they believe.
					\item And you can do this for spiritual experience and religion, too.
					\item And this leads to an argument about your spiritual experiences, about your experiences with the Holy Ghost, that is really appealing to people, especially if you've been to graduate school, because there you tend to intellectualize everything.
					\item The argument is basically, when we put people in a brain scanner and they say that they're ``feeling the spirit,'' we can see that different areas of their brains are active during those experiences of feeling.
					\item It's a way of explaining away your spiritual experiences as just a product of your brain, or a product of evolution, or something like that. And if those are the ways you tend to explain other things, then this argument will seem very powerful to you.
					\item And it will also call into question everything else: your experience with the atonement, repentance, the temple, and so on.
					\item It was powerful to me, for a long time. When I first started thinking about it, it made me go back to all those mission experiences that I was mentioning earlier, and think about how maybe all of that was just me, and not the spirit. All of those words were just coming from me and not some external source.
					\item This is basically the going explanation that people in philosophy, psychology, neuroscience, and so on, give for spiritual experiences.
					\item I've come to see this argument a little differently.
					\item (How I think about it---if God did exist, and did want to communicate with me, why wouldn't he be using my brain? How else would he do it?)
				\end{compactitem}
				
			\item So I hope this answers a little bit about why I believe---I was saying before that I've had experiences with the Holy Ghost, and they've been very powerful.
				\begin{compactitem}
					\item There are lots of alternative explanations for these experiences, like anxiety, or my brain, or evolution, or something like that.
					\item But I don't find these alternatives very convincing.
					\item The more I learn about the science side of these things, the more I feel that these alternatives are reaching pretty far and ask you to make some pretty big leaps.
					\item Maybe in the future these alternatives will be more interesting because we'll know more, but they aren't that interesting now.
					\item There are limits to what you can explain by using evolution and neuroscience, and those ways of thinking are not as powerful as they appear when it comes to religion and spiritual experience.
				\end{compactitem}
				
			\item The second part of my answer was about answering questions---the gospel has answered questions for me.
				\begin{compactitem}
					\item Let me tell you a bit more about this.
					\item For me, reading the scriptures, and reading from people like Joseph Smith and Brigham Young, is exciting. I feel the same excitement with those things that I do with other things that I love to study.
					\item Philosophy hasn't taken away from this. I think it's only added to it.
					\item (Philosophy is methods and ways of approaching problems.)
					\item (The fact that the majority of philosophers don't believe in God, or the majority of neuroscientists don't believe in God or a soul or whatever, doesn't really tell you very much.)
					\item (How philosophy teaches you to ask questions, but also to figure out your assumptions so that you can ask the \textit{right} questions.) 
					\item (This has led me to understand a lot of things differently.)
					\item (Example of the story about finding the BY quote.)
					\item BY quote: ``When God speaks to the people, he does it in a manner to suit their circumstances and capacities...Should the Lord Almighty send an angel to re-write the Bible, it would in many places be very different from what it now is. \textbf{And I will even venture to say that if the Book of Mormon were now to be re-written, in many instances it would materially differ from the present translation.} According as people are willing to receive the things of God, so the heavens send forth their blessings. If the people are stiff-necked, the Lord can tell them but little.''
					\item (The gospel is just this fountain of truth that rewards me whenever I ask questions.)
				\end{compactitem}
				
			\item And this is another reason why I still believe: I have always thought that the true gospel of Jesus Christ would reward me for asking questions.
				\begin{compactitem}
					\item And that is what the restored gospel does.
					\item It rewards you for asking questions, by giving you insight and understanding and clarity.
					\item And this has been a very powerful source of inspiration for my faith.
				\end{compactitem}
				
			\item What is the point of being a person of faith? 
				\begin{compactitem}
					\item (Experiences, questions, and knowledge that you cannot get any other way.)
					\item There is no substitute for these things. If you don't have them, you are missing out on an extremely significant part of human life.
				\end{compactitem}
			
			\item My beliefs are still changing and growing. And this to me is another reason why I still believe: because I feel my knowledge and experience expanding.
				\begin{compactitem}
					\item Early in his ministry, Christ’s disciples come to him and ask about fasting. John’s followers fast, they say, and the Pharisees fast too, and we know fasting is important---but you aren’t commanding us to fast, Jesus! Why not? Why don’t you say anything about fasting, when other important people seem to?
					\item In response Christ makes a remark about how he’s currently present with them, but to make a larger point he then also uses an example of wine. Fresh wine is still fermenting and so releases a lot of gas; if you put it into a wine-skin or other container and try to stop it up, the gas will get trapped. If it’s a new container then that won’t be a problem, because the container can stretch along with the wine. But if it’s an older container and the material has already stretched to its limit, then you’re going to have trouble, because the wine could release so much gas that the container will burst.
					\item Christ’s point is that new ways of thinking sometimes cause that same trouble for old ways of doing things. The new doesn’t always fit with the old, and might even make us uncomfortable---the Spirit can tell us new things, uncomfortable things, even disturbing things about ourselves or who we really are. Some of that may not sit nicely with what we already believe or the ways we’re already quite comfortable living. But sometimes we have to be shaken a bit roughly out of old patterns, either because we’re so deeply entrenched in them that it’s difficult to get out, or because the patterns are so familiar to us that we don’t even realize they’re there.
					\item Joseph Smith made some comments in this spirit near the end of his life. On 13 August 1843 he said that ``It is the constitutional disposition of mankind to set up stakes \& set bou[n]ds to the works and ways of the Almighty.'' The implication there is that the tendency to set up ``stakes'' or limits to what God does is bad, but two weeks later he went farther than implication and just said it: ``To all those who are dis-posd to say to set up stakes for the almighty--will come short of the glory of god'' (27 August 1843).
					\item Why would we come short of the glory of God---because God is punishing us, or something like that? Nope. It’s because when we ``set up stakes'' we’re no longer willing to accept the new, and when we aren’t willing to do that, there’s no transformation or change left for us. There’s only stasis, which in the long run isn’t getting us anywhere.
					\item One of the major challenges of mortality is figuring out how we can make ourselves as much of a new wine container as possible, as often as possible, to stretch to receive what we’re given. It’s frustrating and annoying, especially as you get older, and not easy. But a bit at a time we can do it, and one of the primary functions of the atonement is to induce in us the childlike malleability of character that will allow us not just to receive the new but become the new: the new man, the new woman, the new creatures, the new life.
				\end{compactitem}
				
			\item Christ is real.
				\begin{compactitem}
					\item The atonement is real.
					\item Covenants are real.
					\item We can experience the reality of these things for ourselves, and we can trust that what we experience is real and nothing is deceiving us.
					\item And that is why I believe.
				\end{compactitem}
				
			\item Let me start out by saying that, w
				\begin{compactitem}
					\item (Give background to myself.)
					\item Lots of experiences with the spirit on my mission.
				\end{compactitem}
			
			\item Let me introduce myself a bit more since not all of you have had the misfortune of being in one of my classes.
				\begin{compactitem}
					\item I was born in Sandy, Utah, and when I was 13 we moved to Vernal, which is a small oil field town on the eastern side of Utah.
					\item We lived in the more residential part of the town, but even that meant that we had a field full of sheep across the street from our house, and we used to have home run derbies into their field with bonus points for bonking a sheep.
					\item I went to Uintah High School and played lots of sports but also played lots of video games and read a lot of science fiction.
					\item I still love all of those things and I extremely occasionally stream on Twitch, but I very frequently play the Zelda Tears of the Kingdom game with my kids, and I'm trying to teach myself how to write fiction but I'm terrible at it.
					\item Just a few weeks after I graduated high school I left for my mission in Chile.
					\item More about that in a moment.
						% \begin{compactitem}
							% \item Back then you had to get special permission to go when you were 18 but I got it and left.
							% \item I look back and am not really sure why I served. I'm glad I did, but neither of my parents did, and they never pressured me and we actually almost never even spoke about it.
							% \item I guess I just did it because it was the next thing to do.
						% \end{compactitem}
					\item I got home in 2007 and started at BYU just a few weeks later.
					\item I thought I wanted to be a philosophy major and took courses in philosophy, and you'd think that I would have taken any philosophy course I could but that isn't what happened.
					\item I took the bare minimum number of required classes for the major and I graduated without any minor at all.
					\item But I graduated with 221.5 credits---I had to go back and check my transcript, and I remembered right, I had 221.5.
					\item 221.5 credits is almost double the number of credits you need for an entire degree, and it's because I just loved taking classes in all sorts of areas. 
					\item But what I loved the most was \textit{languages}.
					\item I had to check my transcript for this took, but I took 81 credits of language courses in my four years at BYU.
					\item I took courses in:
						\begin{compactitem}
							\item Latin
							\item Greek
							\item Spanish
							\item German
							\item French
							\item Icelandic
							\item Hebrew
							\item Aramaic
						\end{compactitem}
					\item I loved every minute of those courses, and I still love them.
					\item I can read Greek and Latin well and Hebrew well enough to understand what is going on and use different dictionaries and sources.
					\item Reading and thinking about the scriptures in their original languages has had a huge impact on my thinking about the gospel and has also just been really fun.
					\item Studying ancient languages is difficult, but it pays you off. No effort is wasted, especially when you're able to look at something like the New Testament when you're done.
					\item Anyway, during this time I also met Alicia, who was in my ward, and a year later we started dating and were married in 2010. 
						\begin{compactitem}
							\item There was just instant chemistry between us---interesting conversation from the very first time we spoke, physical attraction, we had fun together, just enjoyed being around each other.
							\item And that's still true today. To this date my wife is the only person I have ever met that I don't get tired of being around.
							\item We have four kids.
						\end{compactitem}
					\item For some reason I felt like I wanted to study philosophy in graduate school, which was weird because I wasn't that great at it.
						\begin{compactitem}
							\item I was much better at languages, and so it seems like going into a classics PhD, which is studying ancient Greek and Latin, would have been a better fit.
							\item But philosophy always interested me and it felt right to keep going.
						\end{compactitem}
					\item I applied to graduate programs and did very badly. So I ended up going to a masters program at Tufts, which is a school near Boston, some of you might know Jaden Nelson, who's there right now for dental school I think.
					\item I went there because my wife is from that area so we saved money by living with her parents and it was easier for her to get a job because she knew a lot of people there.
					\item We lived with my parents in law for the first year, and then moved to an apartment down the road that was a dump, but we were having our first baby and felt like we needed to be out of there.
					\item But I kept studying philosophy and I started to take it seriously and really try to understand things.
					\item I finally got a clue and applied to other PhD programs again and got in and so we moved to North Carolina, where I did my PhD at \textit{the} university of North Carolina.
					\item And by that I mean of course, Duke University.
					\item I spent most of my time studying neuroscience and psychology, and also the history of science, which is probably my favorite subject ever.
					\item I loved being at Duke, I loved living in North Carolina, and I loved the three most important things about the south, which are barbecue, fried chicken, and biscuits.
					\item I finished studying at Duke in 2019 and then came up here that summer and started here fall of 2019, which meant I had one complete semester before the pandemic hit.
					\item There are still a few people around who were in my classes pre-pandemic.
				\end{compactitem}
				
			\item So when I asked some students what sorts of things they might like to hear about, one person said, you should talk about why you still believe, even though you're a philosopher.
				\begin{compactitem}
					\item And I get what they mean---when I told people in my old home ward in Vernal that I wanted to study philosophy in graduate school, one woman, she had been my junior high typing teacher, told me, ``Oh, don't do that. You'll go inactive.''
					\item And I was like, ``Uh, gee, thanks for the advice.''
					\item And I thought, I don't think I will, but I guess I don't know. I can't predict the future.
					\item So it does feel a little bit like I'm addressing the topic not of why I believe, but why I \textit{still} believe---why would I continue to believe, after having studied the things I have?
					\item I'll try to talk about that soon.
					\item But the ``still believe'' thing is relevant in another way, because as I look around, I would say between 30 and 40\% of my college roommates, college friends, and extended family who are around my age---so basically my peers over the sixteen years since I got back from my mission---around a third of them or so have chosen to no longer affiliate themselves with the LDS faith.
					\item That number isn't an exact count, but it's probably about right, and I expect it will continue to rise.
					\item So when I think about the topic of ``why I believe,'' what I also end up thinking about is, why do so many people I know no longer believe? It has nothing to do with graduate school or professional field, either, it's just a general thing.
					\item Why does there seem to be so few of us left standing? Of the people I'm thinking about when I say this, the men all served missions and so did many of the women, and all of them were married in the temple.
					\item So what happened? What is happening to my generation, and why does it keep happening? Why is it that every like four or five months I hear about an old college friend or somebody we were friends with in a previous ward, and they're announcing on Instagram that they no longer consider themselves LDS?
					\item It's very sad for me, because you can't do anything about it. You're just a witness to this massive car accident happening in slow motion, leaving all this hurt and pain and bitterness, and you have this horrible feeling that there's just something that has gone very, very wrong.
					\item And it feels like a war of attrition. Sometimes it feels like you're just waiting around until your number gets called and you're the next one up, it's your turn to make the social media post and write the email to the family and then start making choices about how you're going to live your life when you no longer need to consider the interview questions for a temple recommend.
					\item So when I consider the topic of why I believe, or why I still believe, all of this is in my head, and I'm thinking not only about myself but about all the amazing and wonderful people I've known over the years.
				\end{compactitem}
				
			\item So to actually answer these questions, the first thing I want to say is that philosophy and science are just \textit{methods} or \textit{ways of studying things}, really. They tend to focus on certain topics, but at the core they're just ways of looking at things.
				\begin{compactitem}
					\item You can study anything you want with philosophical methods, and you can study anything you want with scientific methods.
					\item We use these methods to study certain things. The methods fit certain areas better than they fit others.
					\item It's true that when people go into a certain field, they tend to adopt many of the views and ideas of people in those fields.
					\item Sometimes that's because those views and ideas have been arrived at by the best methods and so they are pretty well-tested, but that isn't always true.
					\item Groups of people just tend to kind of even out their opinions and beliefs over time, and that process isn't necessarily a rational one.
					\item In other words, the fact that most philosophers don't believe in God doesn't really tell you anything about whether that opinion is the correct one. It just means that most of the people in the group don't have that belief.
					\item The same is true of neuroscience. I don't know what the percentages are, but even if a majority of neuroscientists didn't believe in God, this wouldn't really tell us very much.
					\item And in fact, I'd bet that most neuroscientists haven't ever given the issue of God's existence any serious thought, so there's very little reason to care about what their opinions are.
					\item There are exceptions---I had a professor at Tufts named Dan Dennett, who is a world-famous atheist who has written many books on religion and religious belief.
						\begin{compactitem}
							\item He knew I was religious and LDS, but he never brought it up.
							\item We talked about religion a lot, and we always did it from this kind of detached perspective of an observer.
							\item It was my dad that was worried about me this time. He was worried I'd become a little Dan Dennett and become an atheist.
							\item I did learn a lot from him and he's actually a huge influence on my thinking even today, but not in his atheism.
						\end{compactitem}
					% \item On the other hand, I do think that both philosophy and science have been extremely important to me and my spiritual life, mainly because the things I've studied have given me a lot of tools for thinking through the different questions that I've had or that other people have asked me.
					% \item The fact that there are these questions and problems to solve isn't very remarkable.
					% \item In every field of study and in every area of human experience there are unsolved problems and questions, so it isn't a surprise that we encounter those things when thinking about our religious beliefs.
					\item The way I see philosophy and science is more like, hey, look at all these new ways I have of understanding my beliefs! This is great! I'll be able to learn so much and understand so many new things!
				\end{compactitem}
				
			\item And I've tried to take that attitude, like an eagerness to learn into both my professional life and my study of the restored gospel.
				\begin{compactitem}
					\item When it comes to the sorts of questions and problems that LDS people ask about their faith and their church, my attitude has always been to just jump in head first.
					\item I've never seen any reason not to. Why not try to answer questions? Why not try to figure things out? 
					\item I remember I'd read what my mom would call ``anti-Mormon'' stuff before my mission, and my mom would say, why do you read that stuff?
					\item And I'd think, well I have to know what the counterarguments are. I have to know what the problems are, or else how can I solve them? I can't think about them if I don't know what they are.
					\item Story about Jehovah's Witness bible on my mission
					\item This has just always seemed like the right thing to do, and that attitude has been reinforced by studying philosophy and science, because nothing is really off-limits for them.
					\item You can ask any question, you can examine anything.
					\item Not everything will necessarily be worth the time, but you could look at anything.
					\item So that's always been my attitude. %even though it wasn't my family's attitude.
						% \begin{compactitem}
							% \item In my family we didn't talk about much of anything, ever.
							% \item It's not this way now, but back then we all loved each other but we really kept a lot of things to ourselves about most things.
							% \item My wife's family is really different. She grew up in a family where people talked openly about almost everything, so for her this attitude is a bit more natural.
							% \item I guess for me it was a bit more learned, it took more to get it started.
							% \item If you're more like me and my family when I was a kid, that's fine. There's nothing wrong with being that way. And there's nothing wrong with being like my wife's family either. 
							% \item If you like to talk and ask questions and explore, it may feel like it's hard to find people you can feel comfortable doing that with.
							% \item But you will eventually find people.
							% \item And you can always come and talk to me, since everyone here is welcome to do that.
						% \end{compactitem}
				\end{compactitem}
				
			\item Anyway, like I said, whether it's because of my profession or my personality or whatever, my attitude has just been to get into the different questions I had and other people had about the church.
				\begin{compactitem}
					\item I mean like stuff with church history, social issues, theological stuff, scriptures, Book of Mormon, whatever, all of those things.
					\item I just started reading about them and tried to understand the problems. I tried to actually think about them.
					\item I see now that I've taken a pretty active approach. That may not necessarily work for everyone.
					\item My wife's approach is different. She's like, I'm going to write this question down, and put it in this jar. And then I'll just trust that someday that question is going to be answered.
					\item I don't think there's anything wrong with doing that, it's just not something that would work for me.
					\item When it came to questions about the church, or faith, or whatever other problems people could have with religious belief, I have always thought about it like this---
					\item I was like, we say this is true. We say this is the true church, we say the atonement is real, we say this is true.
					\item And so I think, well, if it \textit{is} true, then it \textit{better} be able to answer my questions, because if it \textit{can't}, then that's a problem.
					\item That doesn't mean I can expect a detailed answer to any possible question---no one has a theory of everything.
					\item But my mindset was, if I seriously investigate these questions about the restored gospel, and I sincerely try to understand the issues, then I better get something out of that.
					\item I better walk away with some real insight and some real understanding, because \textit{that} is what the true gospel would give me.
					\item The true gospel of our Lord and Savior Jesus Christ would \textit{reward} me for asking questions.
					\item Think about that---it would \textit{reward} me for asking questions, and the reward would be new \textit{insights} and new \textit{understanding} and a new \textit{clarity of vision} about things that \textit{really matter}.
					\item And the reason why it would reward me this way is that, if the gospel is true in the way we say it is, then it must be the right way to see and think about things.
					\item And the right way to see and think about things will always, always reward careful questioning, because it's just a fountain of truth, and it isn't ever going to run dry.
					\item Example: Story about finding the BY quote
						\begin{compactitem}
							\item You see this all the time in the history of science.
							\item One way you see it is in how people used to think about \textit{motion}---I mean like how objects actually move.
							\item Unfortunately we don't have time to get into this here, but motion is an extremely puzzling phenomenon once you start to look at it.
							\item Philosophers before Socrates thought about it, Aristotle did, astronomers like Ptolemy thought about it, medieval philosophers thought about it, like everybody did.
							\item And it stumped all of them because they carried around these incorrect assumptions about it for such a long time, and they just had no idea that their assupmtions were wrong.
							\item They were super smart people, but their starting-points were just all messed up.
							\item It isn't until people have been thinking about it for thousands of years that people like Copernicus, Kepler, Galileo, and Descartes start to finally begin asking the \textit{right} questions.
							\item People have been wondering about motion for ages, but they've been asking all the wrong questions, because of the incorrect assumptions they had about the universe to begin with.
							\item And often what happens when you ask questions that are based on incorrect assumptions is that your questions and your attempts to answer them turn out to be massive red herrings that lead you down all sorts of wrong paths.
							\item But if you are able to \textit{fix} those assumptions and start asking the \textit{right} questions, you get results almost immediately.
							\item You're not going to believe this, but one of the first people to ask the really important, really clear questions about motion was Descartes.
							\item Everyone here thinks he's an idiot because of all the doubt stuff you read about in Reason and the Self, but he actually did way more than that.
							\item Descartes finally asks the questions correctly, and then do you know what happens after that?
							\item In less than fifty years you get Isaac Newton, maybe the most important scientist who ever lived.
							\item And by building on what Descartes and others had done, Newton asked the questions that were exactly the correct questions that people should have been asking for over two thousand years, but they couldn't ask those questions. I mean they literally could not because those questions could not even be asked within the framework of the assumptions they had.
							\item Over time people fixed the framework and fixed the assumptions, and what do Newton's questions give you?
							\item In 1687 they give you the laws of motion, which are principles so fundamental and so broad that they apply to literally \textit{everything} in the entire \textit{universe}.
						\end{compactitem}
					\item So there's an extremely important pattern here for us to learn.
						\begin{compactitem}
							\item Very frequently, the questions we are asking about not only the gospel but about anything at all are the \textit{wrong questions}.
							\item And they're not wrong because we're stupid or can't handle the truth or whatever.
							\item They're wrong because they are based on incorrect assumptions we make about the gospel, about ourselves, about the world, about science, or about lots of other things.
							\item Our questions are sincere and well-intentioned but they turn out to be red herrings, just like there were so many red herrings in the history of the way people thought about motion.
							\item And so the next step in the pattern is, we sometimes are able to learn things that help us fix our assumptions. We sometimes are able to think through things, or talk to another person, or learn something crucial that fixes some or all of our assumptions.
							\item (When you think about it, this is most of what therapy is---it's trying to get to these harmful assumptions about ourselves and fix them.)
							\item Then, as the pattern goes, we become aware of our assumptions and we fix them, and now, we are finally in a position to ask the \textit{right} questions.
							\item And those won't be red herrings because they're built on correct assumptions about ourselves, the gospel, the world, and so on.
							\item And just as happened with Descartes and Newton, whose questions led to world-shattering insights into the very nature of the universe, that is what the right questions about the restored gospel can do.
							\item They can lead to life-altering discoveries and world-shattering insights.
							\item And when you think about it, it just has to be this way. It couldn't be any other way.
							\item The gospel is as broad as the universe and as deep; if it's true, its truth touches the foundations of everything that exists.
							\item So we ought to be able to see this same pattern play out: if I take the questions seriously, and I am diligent in trying to understand things and ask the correct questions, then I have a right to expect the same world-shattering insights from the true gospel.
							\item And that is what the restored gospel of Jesus Christ gives. That is what it offers us.
							\item Remember, the true gospel will reward you for asking questions, because it will be able to answer them with clarity and insight.
							\item The true gospel is a fountain of truth which will never run dry---it's like the position Newton was in at the end of the seventeenth century.
							\item All of a sudden everything clicks into place, and he's given just astonishing intellectual power and penetrating discernment into the universe.
							\item The gospel of Jesus Christ is like that. If it's true, it has to be like that. 
						\end{compactitem}
				\end{compactitem}
				
			\item And this, for me, has been the great contribution of philosophy and science to my spiritual life and my faith and my religious belief: they have given me the tools and ideas that I needed in order to figure out what was mistaken about the quetions I was asking, so that I could start to the right ones instead.
				\begin{compactitem}
					\item This is why studying philosophy and science and the history of science and other stuff has only strengthened my beliefs, rather than weakened them or hollowed them out.
					\item What has happened is that what I have learned professionally has ended up giving me the tools that I needed to both ask and answer serious questions about the gospel.
					\item One of the things that philosophy teaches you is what is actually involved in asking the \textit{right} questions, and then answering them \textit{deeply} and \textit{systematically}.
					\item It teaches you to unearth assumptions that are very deeply embedded in how you think about things and in how you see the world, and then teaches you to test those assumptions and either reject or change them so that they're better.
					\item And once you fix those assumptions, it's like going from being Aristotle to being Descartes and Newton, like I said. You are finally, at long last, in a position to ask the right questions.
					\item So philosophy doesn't just teach you to question things, or if it does, it ought to be doing more than that. What it should be doing is teaching you how to ask the \textit{right} questions, and teaching you how to figure out which ones are the right ones in the first place.
					\item Those will be the ones that will be most likely to reward the effort you spend in asking and answering them.
					\item But the truth is that most people have no idea at all what it would take to carefully consider and investigate their sincere questions about their faith.
					\item And like I said before, it isn't because they're dumb, it's just because no one has ever shown them what that would actually involve. 
					\item The idea that most people have of investigating something, or of really looking into something faith-related, is like watching a TED talk.
					\item If you haven't seen them, TED talks are these video presentations on Youtube that last between five and twenty minutes and they go over a specific topic in science or politics or whatever. Everything is explained at a very general level and it's all made to seem just so nice and pleasant with a little bow on the top at the end.
					\item But my friends, I am here to tell you that you do not belong to a TED gospel; you do not belong to a TED-level faith.
					\item The true gospel of Jesus Christ could never be a TED-level anything, because it's supposed to be the truth of \textit{everything}.
					\item Why would it be easy to get answers from something as broad and as deep as the entire universe?
					\item We're just so used to our systems of thought and ways of thinking being like vending machines, where you just go up to it and punch in a code for whatever your question is and it spits out like a bag of chips or a candy bar for an answer.
					\item But in order for the process to work like that, you have to dumb everything down so much that the answer you end up getting is as satisfying as a bag of chips or a candy bar. You eat those things and enjoy them for like two minutes and then you feel empty and even worse than before.
					\item If you are asking sincere and meaningful and difficult questions, you cannot expect your answers to come bundled in familiar, easy language, totally free of any implications that might challenge the way you normally think about things.
					\item That is how vending machines work, and that is how superficial, TED-level understandings of something work. But that is not how the truth works.
					\item If investigating the truth were that easy, then it wouldn't have taken thousands of years to get to the laws of motion. Aristotle would have found them a long time ago and there would have been no need for Copernicus or Descartes or anyone else to fix all the incorrect assumptions that people had.
					\item But again, the truth doesn't work like that. 
					\item And the restored gospel of Jesus Christ doesn't work like that either, because it is the truth.
				\end{compactitem}
						
			\item So some of the students I was talking to about this devotional were also wondering, what is the virtue of faith? What is the point of it? What is the purpose of being a person of faith?
				\begin{compactitem}
					\item This is a really good question, and actually this is one of the questions that gets more pressing as you get older, because the older you get, the more you start to feel like being a person of faith just doesn't make a real difference to your life.
					\item You can be a good person with or without faith, so why make the effort? What's the point?
					\item I think this is one of the great rewards of faith: being a person of faith puts you in a position to ask extremely deep and meaningful \textit{questions} that other people are simply not able to ask.
					\item It isn't that they don't \textit{want} to ask them; it's that they \textit{can't}.
					\item They cannot ask them because they are blind to the reasons why you would ever need to ask them in the first place, and in some cases they don't even have the l,anguage to frame the questions at all.
					\item And likewise, being a person of faith and knowing the true gospel of Jesus Christ then puts us in a position to \textit{answer} those questions---not with vague, unsatisfying language, but with actual \textit{principles} that can re-frame the way you see and experience yourself, other people, and the entire world.
					\item As a result, there will be knowledge available to you, there will be truth available to you, there will be experiences available to people of faith that are simply not available to other people.
					\item They will never ask those questions and they will never have those experiences.
				\end{compactitem}
				
			\item This is also the reason why I think religion and spirituality, on the one hand, and science and whatever else you want to put with it, on the other hand, are not in conflict with each other. They aren't in conflict and cannot substitute for each other because they are about different domains of human experience. 
				Professor Lee and I taught a course on science and religion last semester and we ran into this problem a lot.
				People think, you need religion for morality, you need religion for community, and 
				atheists know this and so write books called things like "the atheist's guide to reality" or whatever.
				people have been
				
			\item This is exactly one of the things that philosophy does best, and teaches you to do. It teaches you not just to identify your assumptions but to try to fix them, and it gives you the tools to do that.
					And for me this has been a very empowering process. I feel like I've been involved in this process for years and expect to be involved in it for as long as I live, because we're constantly learning and gaining new information and we constantly need to revise what we think about things in response to what we learn. This is as true for the gospel as it is for anything else.
					And I think there's some risk in not doing this, because many people, later in life, will have to confront this, and the confrontation will take place whether they know it's happening or not.
					In many cases it happens like this.
						You spend [obstacle course]
						You finish and finally you can rest, and you start to pay attention to other stuff.
						And one of the things you start to pay attention to is your experience of your faith.
						What happens to a lot of people is that you experience something uncomfortable about your faith, and you call it doubt because you don't have a good term for it and becaus we often oppose faith and doubt. But in many cases it isn't doubt. It's friction between your experience of the world and the things you believe about the world, including things connected to the gospel.
				
				\item It's almost like we don't take the gospel seriously ENOUGH. We don't really consider what it would be like if the fundamental claims of the gospel, like about the atonement and other things, were actually TRUE, like really and honestly TRUE, and not just stories. We talk about it all like it's just stories sometimes, like the pre-existence and the celestial kingdom and the atonement and everything is just stories, like a set of mythological tales that make us feel good when we go to bed at night or something.
				
				\item But like we've been saying, that's not what the gospel of Jesus Christ is. It is not just a set of nice stories that your parents or the missionaries told you.
				
				\item The gospel the great underlying and unifying truth of the universe. And we have to remember to treat it that way. If you take the gospel medium-seriously, you'll feel the Spirit sometimes and you'll read the Come Follow Me assignment and so on. But if you take it totally seriously---if you take the fundamental claims of the gospel as just straight-up true, as things that either have happened or are real---then you will begin to glimpse that great underlying and unifying truth. You will begin to grasp the clarity and the vision and the beauty and the power of these truths.
				
		\end{compactitem}
		
		\begin{compactitem}
				
			\item The topic of this devotional series is ``Why I Believe.''
				\begin{compactitem}
					\item That question I can answer really easily. 
					\item Christ is a real person, and the atonement is a real event.
					\item And I believe in those things because they're real, just like I believe in the existence of, oh, I don't know, corn dogs, because they're real too.
					\item And if I ended right now this would be a talk that was really well suited to my own attention span.
				\end{compactitem}
				
			\item But when I think a bit more about the topic, it's actually a difficult one for me, because I can't separate it from the question of why so many people I know \textit{no longer} believe.
				\begin{compactitem}
					\item So I guess I want to talk about that.
					\item I have written and re-written this a few times, though, so I've had a hard time deciding what to say.
					\item When I took this job a few years ago, I knew that I would sometimes have these sorts of discussions with students---not necessarily like this, at a pulpit, but in class or in my office or whatever.
					\item And the one thing I've always tried to do with students is be totally honest. So I'm going to do that today. I don't know how uplifting this is going to be, but it will be completely honest, and I hope you'll appreciate that, even if you appreciate nothing else about what I say.
					\item So---I turned 37 a few weeks ago.
					\item I started my mission a few weeks after I graduated high school in 2005, and then got home in 2007 right before fall semester started.
					\item That means it's now been just over 16 years since I finished my mission.
					\item In that time, I went to college, got married, graduated, went to graduate school, and we had four kids.
					\item But also in that time, between 30 and 40\% of my college roommates, college friends, and extended family who are around my age---so basically my peers over those sixteen years---around a third of them or so have chosen to no longer affiliate themselves with the LDS faith.
					\item That number isn't an exact count, but it's probably about right, and I expect it will continue to rise.
					\item I'm not going to talk about the reasons why these people have made this choice. Some have reasons that are good, others have reasons that are less good.
					\item But when I think about the topic of ``why I believe,'' what I end up thinking about is, why do so many people I know no longer believe? Why does there seem to be so few of us left standing? Of the people I'm thinking about when I say this, the men all served missions and so did many of the women, and all of them were married in the temple. So what happened? What is happening to my generation, and why does it keep happening? Why is it that every like four or five months I hear about an old college friend or somebody we were friends with in a previous ward, and they're announcing on Instagram that they no longer consider themselves LDS?
					\item It's very sad for me, because you can't do anything about it. You're just a witness to this massive car accident happening in slow motion, leaving all this hurt and pain and bitterness, and you have this feeling that there's just this horrible feeling that something has gone very, very wrong.
					\item And it feels like a war of attrition. Sometimes it feels like you're just waiting around until your number gets called and you're the next one up, it's your turn to make the social media post and write the email to the family and then start making choices about how you're going to live your life when you no longer need to consider the interview questions for a temple recommend.
					\item It's just so frustrating, because fifteen or twenty years ago, we were the ones being called the ``rising generation,'' ``the ones prepared for these latter days,'' and all that other stuff that you guys get called now.
					\item And it doesn't feel like we are that or ever were that.
					\item Now in saying this, I don't want to give the impression that I consider myself better or more righteous than those who have taken different paths.
					\item In fact I know I'm not better, and this is one of the most frustrating parts, that personal righteousness doesn't seem to be a factor in what I'm talking about here.
					\item It's like it doesn't protect you.
					\item I'd love to say, ideas protect you, good ideas protect you, and sometimes I believe that and sometimes I don't.
					\item The truth is that the explanation for why I still believe has less to do with deliberate choice and more to do with circumstances and things I happened to learn at different times,
					\item And that's a little scary to think about, that my continued belief has less to do with deliberate choice and more to do with circumstances and things that just happened to occur.
					\item But as I look back, at least now, that's what seems to be true.
					\item Maybe, probably even, the Lord has provided these opportunities for me to learn. But I don't feel uniquely deserving of them and I'm not sure why he hasn't done the same for other people.
				\end{compactitem}
				
			\item So the simplest answer to the question of why I still believe is that, by either a stroke of good fortune or some very convoluted divine intervention, I ended up studying in my profession basically the exact ideas that I would later use to understand my own spiritual experience as well as the restoration itself.
			
			\item neuro
			
			\item history of science
				examples
				
			\item changing views of the gospel (BY quote, other sutff)
				
			\item This is going to sound really strange, but the biggest part of why I still believe next major reason that I still believe is that, one day, completely by random chance, I came across a quotation from Brigham Young that changed the way I thought about the gospel and put my thinking about it on a completely different track.
				\begin{compactitem}
					\item We got married and lived in this tiny studio apartment in Provo, Utah.
					\item Our bed was in a loft above the kitchen so every morning I would go downstairs and sit at the table and read the scriptures.
					\item I had this old desktop computer with a database on it that could search through a lot of sources, including the Journal of Discourses.
					\item If you don't know them, the Journal of Discourses is a set of books that have general conference talks and other talks from church leaders given from about the 1850's until the 1880's or so, around there.
					\item Anyway, one day I was just using the database to search sort of random stuff, just playing around. I don't remember what I was searching for in particular.
					\item But the search brought up a remark Brigham Young once made about the Book of Mormon and what it would have been like if it had been re-translated later.
					\item The remark seemed interesting, but I didn't write it down or anything, I just moved on.
					\item But something about it stuck with me, and I thought about it a couple months later, only when I went to find it again, I couldn't find it. It took me a long time of re-searching through that database to find it.
					\item Finally I did, and this time I wrote it down so that I would have the reference in the future.
					\item I'm going to read that remark to you now, and then tell you a little about it. It's from a talk Young gave in 1862, and the source is Journal of Discourses volume 9, page 311:
						\begin{compactitem}
							\item ``When God speaks to the people, he does it in a manner to suit their circumstances and capacities...Should the Lord Almighty send an angel to re-write the Bible, it would in many places be very different from what it now is. \textbf{And I will even venture to say that if the Book of Mormon were now to be re-written, in many instances it would materially differ from the present translation.} According as people are willing to receive the things of God, so the heavens send forth their blessings. If the people are stiff-necked, the Lord can tell them but little.''
							\item Let me read the part about the Book of Mormon again---remember that he said this in 1862, 32 years after the book was published.
							\item (Read bold part again.)
							\item That idea absolutely blew my mind and all of a sudden I had like a million questions---like how could the translation be different? Wasn't it written on plates? How could anything be different?
							\item That idea started to work on me, and it worked and worked and it hasn't ever stopped.
							\item Eventually it caused me to see the scriptures differently and to see many events in the history of the restoration of the gospel differently as well.
							\item I believe that the reason I've never had anything like a testimony or faith crisis is that I found this idea, and other ideas like it, and they and the Spirit helped me to think about things in new ways such that nothing like a crisis could ever even arise.
						\end{compactitem}
					\item And now, looking back, I'm probably the only person in the church whose faith Brigham Young \textit{saved}, instead of \textit{destroyed}.
					\item I continue to learn a great deal from Brigham Young and Joseph Smith.
					\item This is why I teach classes on them---because I want to try, even if I do it badly, to communicate some of what I've learned from them, and some of the immense power that we find in their thinking about the restoration.
					\item I'm still trying to work out the implications of these ideas about revelation. The scriptures are full of incredibly important ideas about revelation that we very frequently gloss over.
					\item In my opinion, the way to answer difficult and troubling questions about the church is to study \textit{revelation}, and the way it works. For me at least, many puzzling things have dissolved away as I have come to understand the principles by which revelation operates.
				\end{compactitem}
				
% ----------------------------------------------------------
% ----------------------------------------------------------
% ----------------------------------------------------------
% ----------------------------------------------------------
% ----------------------------------------------------------
				
			\item The first important thing that happened to me was that, as a missionary, I started to take the scriptures really, really seriously.
				\begin{compactitem}
					\item I grew up LDS and we read the scriptures as a family pretty consistently. 
					\item And I read them individually, too. I read the Book of Mormon a couple times as a kid, and must have made it through the New Testament at least once too.
					\item I got called to serve a mission in Chile, and I was set to leave just a few weeks after I graduated high school.
						\begin{compactitem}
							\item Looking back, it's hard to even say why I went on a mission. Neither of my parents served and we rarely spoke about it at my home. They didn't pressure me.
							\item It just seemed like the right thing to do.
						\end{compactitem}
					\item I got some stuff in the mail from my mission before I left, asking me to read the Book of Mormon before I got to Chile, focusing on the atonement of Christ.
					\item I did my best to do that.
					\item But when I got down there I met my mission president, a guy named Ken Openshaw, and I went to zone conferences where he would teach from the scriptures.
					\item And in those zone conferences I learned that I knew absolutely nothing about the scriptures, and that I had never really read them at all.
					\item I mean, I had read them, but I had never taken them seriously. I had never taken seriously the fact that they contain revelations given by God through the Holy Ghost.
						\begin{compactitem}
							\item I had never understood that they are the most fundamental way that God gives us information about the plan of salvation.
							\item I had never understood that even the tiniest things in the scriptures, like the presence or absence of a given word, could be extremely significant.
							\item I learned this from my mission president, and out of all the things I gained individually on my mission, this was undoubtedly the most significant.
							\item Let me give you an example, from \hyperref[DC13-1]{Doctrine and Covenants 13}. That section is a report from Joseph Smith's personal history of a messenger who appeared to him and Oliver Cowdery in the early days of the church.
							\item The opening part of the section reads, ``UPON you my fellow servants, in the name of Messiah, I confer the Priesthood of Aaron.''
							\item One time my mission president asked me to hold his scriptures for him and I was flipping through them and came to this section, and I noticed that he had marked the beginning of this section, and noted that it didn't have the word ``the'' in front of ``Messiah.'' 
							\item And he had a note written to himself about whether the absence of the definite article ``the'' indicated that the text is referring to ``Messiah'' as an \textit{office} or \textit{role} that could be occupied by multiple individuals, or whether it's referring to a particular person.
							\item I don't know what the answer to that question is, but I learned that, in order to really study the gospel and the plan of salvation, I had to be paying attention to the scriptures at that level of detail.
							\item And that's really, really difficult, and it takes a lot of effort, and I don't always get there.
							\item But when do, amazing things happen---things just click into place in these moments of absolutely beauty and power when you're just stunned into silence by the grandeur and the expansiveness of the plan of salvation.
						\end{compactitem}
					\item The scriptures are still extremely important to me. I read almost nothing else, at least when it comes to the gospel, though I'll tell you about a few exceptions later.
					\item So to answer the question of why I still believe, the start of the answer is that the scriptures still provide me today with the same excitement and inspiration that they did back when I was a missionary. My interest in them has never diminished.
						\begin{compactitem}
							\item There are many aspects of my participation in church that bore me, but not the scriptures. They have never bored me.
						\end{compactitem}
				\end{compactitem}
				
			\item Not everything that happened on my mission had good long-term consequences. For example, while I was a missionary, I learned to connect powerful emotional experiences with the presence of the Spirit and revelation.
				\begin{compactitem}
					\item So after I got home, I still had that connection in my mind: powerful emotions = presence of the Spirit.
					\item But then in college I developed really severe anxiety. 
						\begin{compactitem}
							\item I had panic attacks, and at times it was so bad I could barely make to class and couldn't pay attention.
						\end{compactitem}
					\item This was really a challenge for trying to discern the Spirit.
						\begin{compactitem}
							\item Anxiety brings with it many very powerful emotions, including a lot fear and other very unsettling feelings.
							\item And previous to this I had connected feelings exactly like those with the Spirit, like with being warned about something or whatever.
							\item So now, my anxiety and my thinking about the Spirit are basically constantly putting me in a position where I think something is spiritually wrong: my feelings are unsettling and uncomfortable, and I assume that's the Spirit.
							\item And I wonder, am I the problem? Am I not righteous enough? Why do I always feel like the Spirit is telling me that something is off?
							\item The culmination of this problem came when I was trying to figure out whether to get married to Alicia, who I am in fact now married to.
							\item I loved being with her, we always had fun and really interesting conversations, we were attracted to each other, everything was there.
							\item But when I prayed about it, all I got was unsettling feelings, and so there was this huge conflict between what I knew I wanted, which was to be with her, and what I felt like the Spirit was telling me, which was that something was off. 
							\item At one point I finally felt like I received a sort of lukewarm ``yes'' from the Spirit and I told Alicia about it and so we were sort of unofficially engaged.
							\item But not that long after I broke off the unofficial engagement because I still felt like my feelings weren't right and the Spirit was warning me about something.
							\item I wanted to be with her, so I was miserable, because it felt like the Lord was telling me I couldn't have something that I both wanted to have and that was clearly a good thing for me to have.
							\item But that wasn't what was actually happening.
							\item I now realize that the Spirit wasn't warning me about anything, and that none of what I was feeling had anything to do with the Spirit.
							\item I didn't know that at the time, and when I say I feel lucky that I made it through our engagement alive, I mean that totally seriously. Everything about my life was a disaster during that time and I regularly thought I was going to actually die---not that I was going to end my life, but that there would be some accident or something that would prevent me from marrying the girl I wanted to marry.
							\item I walked around like that, with those thoughts filling my head, for months. It was horrible and the darkest, most difficult time of my life.
							\item What terrified me the most was just the \textit{future}---how could I know that our relationship would be successful in the long term? How could I know that things would go well for us and that we wouldn't get divorcd?
							\item It's almost like I was expecting the Spirit to guarantee a successful relationship, or to give me proof that we would have that.
							\item But the Spirit doesn't work that way.
							\item If you are hoping that a witness from the Spirit will help you avoid making choices that might be mistakes in the long run, you are hoping in vain, and you will paralyze yourself into inaction.
							\item Do not make the mistake that I did.
							\item We are told to ``study it out in our minds'' (\hyperref[DC9-8]{DC 9:8}) because, in that process, we very frequently discover the information that allows us to make the decision \textit{without} requiring input from the Spirit---so we might discover that one of our options is clearly the best, or that another is totally unworkable.
							\item And in this process we also provide the Spirit with more information and material to work with, so it can add nuances to our understanding that it wouldn't be able to give otherwise.
							\item For me and my choice about marriage, the answer was in my experience: we were compatible and in love. We had fun together and loved being around each other. We wanted the same things. And that was the answer. You can't be more sure than that, and many people that have powerful spiritual experiences related to who they feel they should marry end up divorced or in bad relationships anyway.
						\end{compactitem}
					\item So I learned that for me at least, because of the way I feel things---which is either not at all, and that's most of the time, or in an unbalanced way, and that's the anxiety---I had to sort of sever the connection between my spiritual life and my emotional life. 
					\item To this day I try not to connect strong emotion with the presence of the Spirit because I still find it difficult to sort through my feelings sometimes.
					\item This is another reason that my spiritual life revolves mostly around the scriptures, because reading them is an intellectual rather than an emotional experience.
					\item I think this separation of emotion and the Spirit has had long-term consequences, and the most important one is that when hurtful things related to the church happen, they don't bother me at all. They don't mean a lot to me because the connection is already cut.
						\begin{compactitem}
							\item Please don't misunderstand---I am not suggesting that \textit{you} should do what I did. I think you should \textit{not} do what I did.
							\item But it has changed me. Sometimes hurtful things do happen related to the church and my wife will say, why doesn't this bother you? And I think, it probably should bother me, I recognize why people are upset by it, but it doesn't bother me. I don't feel anything for it.
							\item I guess I had to kill the part of me that would have been bothered in order for the rest to keep going.
							\item I feel like this makes me less human in a way.
						\end{compactitem}
				\end{compactitem}
				
			\item Everything up to this point set me in a particular direction in my thinking about the gospel, since it all happened before I started studying philosophy and science seriously---all of this other stuff was already there.
				\begin{compactitem}
					\item Because the truth is that I did not pay much attention to either philosophy or science when I was an undergraduate.
					% \item I was a philosophy major, but I took the bare minimum number of required credits to complete the major, and I didn't have any minor at all.
					% \item What I focused on instead was \textit{languages}, and I took a \textit{lot} of them.
					% \item I went back to my college transcript to check on this because I was curious.
					% \item I graduated with 221.5 credits, which is almost double the number of credits you needed.
					% \item 81 of those credits were in languages. I took courses in:
						% \begin{compactitem}
							% \item Latin
							% \item Greek
							% \item Spanish
							% \item German
							% \item French
							% \item Icelandic
							% \item Hebrew
							% \item Aramaic
						% \end{compactitem}
					% \item I loved every minute of those courses, and I still love them.
					% \item I can read Greek and Latin really well and Hebrew well enough to understand what is going on and use different sources.
					% \item Reading and thinking about the scriptures in their original languages has had a huge impact on my thinking about the gospel and has also just been really fun.
					% \item Studying ancient languages is difficult, but it pays you off. No effort is wasted, especially when you're able to look at something like the New Testament when you're done.
				\end{compactitem}
				
			\item Anyway, despite taking more than eighty credits in languages in college and just forty in philosophy, I decided to study philosophy in graduate school.
				\begin{compactitem}
					\item This was probably a surprise to people who knew me back then. I wasn't very good at it and had a pretty poor understanding of it, even for someone who had only been an undergraduate major in it.
					\item I applied to graduate programs and did very badly. So I ended up going to Tufts, which is a school near Boston, some of you might know Jaden Nelson, who's there right now for dental school I think.
					\item I went there because my wife is from that area so we saved money by living with her parents and it was easier for her to get a job because she knew a lot of people there.
					\item We lived with my parents in law for the first year, and then moved to an apartment down the road that was a dump, but we were having our first baby and felt like we needed to be out of there.
					\item But I kept studying philosophy and I started to take it seriously and really try to understand things.
					\item The most important thing that happened during this time was that I met a philosopher named Dan Dennett.
						\begin{compactitem}
							\item I had a class about psychology and philosophy of mind with him my very first semester, and loved it.
							\item Dennett is a very outspoken and widely known atheist. He's written a number of books about religion, including one about religious leaders who privately lost their faith but for social reasons they couldn't leave their church.
							\item He's also been called one of the ``Four Horsemen of the New Atheism,'' like, he's one of the major figures in certain atheist movements.
							\item He knew I was LDS, but we rarely talked religion, but when we did, he was always extremely respectful and very kind.
							\item He thought about everything in terms of how we can know things, and in practical terms to, like what difference it would make if \textit{this} idea was the correct one, as opposed to \textit{that} idea.
						\end{compactitem}
				\end{compactitem}
				
			\item Studying neuro
			
			\item Studying history of science
				\begin{compactitem}
					\item Changing concepts, developing concepts, new languages
				\end{compactitem}
				
			\item Revelation, and the essays
			
		\end{compactitem}
		
	% ----------------------
	\subsection{13 August 2023: EQ discussion on the power to heal}
	% ----------------------
		\label{TALKS-13-08-2023}
		
		Based on the talk ``He could heal \textit{me}!'', by Peter Meurs.
		
		\begin{compactitem}
			
			\item Meurs' talk is about the healing power of Jesus Christ, and in particular, Christ's ability to heal aspects of our lives that are not related to sin.
			
			\item Meurs tells a story to illustrate this point, about a time that he was responsible for a very severe accident that hurt a number of different people. Let's read about it.
				\begin{compactitem}
					\item The story starts about two-thirds of the way through in a paragraph beginning, ``In 1990...'' Could someone read that paragraph for us?
					\item We also need to read the next three paragraphs after that one.
					\item So after working with other people and teaching them about the atonement, Meurs came to realize that Christ's power to heal didn't just extend to forgiveness from sin or other sin-related things, but also included the guilt and remorse that he was feeling because of his actions.
					\item We can read the next paragraph in the talk, which begins ``The Savior's healing and redeeming power...''
				\end{compactitem}
				
			\item This is a pretty powerful story of a person's experience with Christ, and how Christ changed that person's feelings about a troublesome event.
				\begin{compactitem}
					\item Meurs' story reminded me of one from the Book of Mormon, where we read about a very similar healing.
					\item We'll go to Alma 24 for this story.
					\item Here, a group of Lamanites are preparing to take up arms against another group who are now calling themselves Anti-Nephi-Lehis.
					\item The king of the Anti-Nephi-Lehis passes away, and gives the kingdom to his son, to whom he ends up giving the name ``Anti-Nephi-Lehi.'' 
					\item The people hold a council together knowing that they are going to be attacked, and they decide that they are not going to fight back, and in fact aren't going to prepare for combat at all.
					\item The king then makes some remarks, which start in \hyperref[ALMA-24-7]{verse 7}.
					\item Let's pick it up in verse 8---could someone read verses 8 and 9?
					\item So they've been convinced of ``sins'' and ``many murders.'' 
					\item Then verses 10 and 11, I'll read those.
					\item (Discuss these verses---hit on God taking away the ``guilt'' from their hearts, and then the idea of ``sufficient repentance'' in verse 11.)
					\item So these people, similar to Meurs, had some very troubling feelings because of actions they had done in their past, and they prayed and Christ's power was essential in taking away those feelings.
				\end{compactitem}
				
			\item I find these stories so interesting in part because they are so foreign to me, in a way. For me, over time, my \textit{feelings} in particular have played less and less a role in my spiritual life.
				\begin{compactitem}
					\item To begin with, I'm not a very emotional personal by nature, I tend not to feel very much. This isn't necessarily a bad or good thing about me, it's just how I am. Maybe some of you are similar.
					\item But one thing I do feel is anxiety, sometimes in extreme amounts.
					\item Tell story of mission---associating spirituality with feeligs of different kinds.
					\item Then getting home and my anxiety blossoming as I connected my immediate feelings with my level of righteousness, and felt like I needed to repent all the time.
					\item Then starting to get panic attacks and, after a long time which included both therapy and medication, realizing that I had to sever a great part of the connection between my feelings, on the one hand, and my spirituality and righteousness, on the other.
					\item Because of my anxiety, I couldn't trust that connection, because one of the most important things to learn about anxiety is that it's misleading---your worries are, at least most of the time, \textit{not} based in reality and therefore are \textit{not} guides for what is real or for your actions.
					\item So now, all these years later, when I start to feel guilt and shame, my first reaction most of the time is to just reject those feelings---not to suppress them, but to just let them sit there in my experience until they go away.
					\item I'm not sure that this is the best thing for me to do, but all this time has shown me that, at least for me, my feeligs are rarely \textit{spirtually significant}---the positive ones are not necessarily indications of the Spirit or something like that, and the negative ones are not necessarily indications of a warning from the Spirit, or of a prompting that I need to repent, or whatever.
					\item So they aren't spiritually significant in that sense.
					\item My question for you all is, \tr{\textbf{How do you know when your feelings \textit{are} spiritually significant? How do you know when you can trust your feelings to be real indications of the Spirit, or the need to repent, or something like that?}}
				\end{compactitem}
				
			\item When we look at the life of the Savior, we see that \textit{feelings} played an extremely important role in the way he lived and in the way he responded to people.
				\begin{compactitem}
					\item In his talk, Meurs discusses some of this.
					\item Roughly halfway through, Meurs is talking about 3 Nephi 17, in a paragraph that begins ``In 3 Nephi 17...'' Let's read that paragraph and the next one.
					\item The key part comes at the end of the second paragraph, that Christ's ``bowels'' are filled with ``compassion.''
					\item Let's go to 3 Nephi 17 in the scriptures to read \hyperref[3NEPHI-17-5]{verses 5--7}.
					\item So Christ sees the people and has a very powerful emotional response to them, and that response compels him to stay, even though he has other things to do.
					\item (Explain the idea of ``bowels'' as the location of powerful emotions.)
					\item In the New Testament, there is a more direct linguistic link between ``bowels'' and this compassionate response.
						\begin{compactitem}
							\item We get the word \gr{σπλάγχνον}, \textit{splanchnon}, which is a noun that basically means the ``inner body parts,'' especially the ones in your stomach and chest.
							\item And then we get the verb \gr{σπλαγχνίζω}, \textit{splanchnizo}, which means ``feel pity, have sympathy for,'' and this is a very common word used to describe Christ's reaction to people who are suffering.
							\item (We actually have a somewhat similar linguistic connection in English, with the words \textit{viscera} and \textit{visceral}.)
							\item Some important locations for the verb are \hyperref[MATTHEW-9-36]{Matthew 9:36}, \hyperref[MARK-8-2]{Mark 8:2}, and \hyperref[LUKE-7-13]{Luke 7:13}.
							\item The word also appears in two extremely important parables in Luke---at \hyperref[LUKE-10-33]{Luke 10:33}, which is the parable of the Good Samaritan; and also at \hyperref[LUKE-15-20]{Luke 15:20}, which is the parable of the Prodigal Son.
							\item We can read a bit from those passages.
						\end{compactitem}
					\item So there is clearly an appropriate emotional response to have to the suffering and difficulties of others---and for someone like me, who's by nature pretty unemotional and has learned to distrust his emotions in some ways, that's a somewhat difficult thing to hear, because it means I have a lot of work to do.
					\item Meurs comments on Christ's response of deep emotion and compassion for others in the next two paragraphs of the talk from where we were, the first one begins, ``I believe that His compassion...'' Let's read those two paragraphs.
					\item You may recognize some of the language from that first paragraph as coming from Alma 7. In fact, there are a number of verses which clearly connect the atonement and this deep, emotional, compassionate response to the difficulties of other people.
					\item Let's take a look at some of these verses and think about them through this lens.
						\begin{compactitem}
							\item \hyperref[ALMA-7-11]{Alma 7:11--13}
							\item \hyperref[ALMA-34-13]{Alma 34:13--16}
							\item \hyperref[MOSIAH-15-7]{Mosiah 15:7--9}
							\item \hyperref[DC45-3]{DC 45:3--5}
						\end{compactitem}
					\item So the atonement, in some way, played a fundamental role in Christ becoming a being who could have this compassionate, emotional response to \textit{every individual}, \textit{in any circumstance}.
					\item This is a state of being that is extremely far from where I am right now, clearly.
					\item Enoch got a taste of this way of feeling and experiencing in \hyperref[MOSES-7-41]{Moses 7:41}.
					\item \hyperref[DC19-15]{DC 19:15--20}
				\end{compactitem}
			
			\item So I take two lessons from Meurs' talk.
				\begin{compactitem}
					\item The first is that the atonement and Christ's power can heal our troublesome feelings, even if it's just a bit at a time over a long period of time.
					\item The second is that part of the transformation we must undergo through the power of the atonement is to become people who can have more of the compassionate, emotional response to other people that Christ has.
						\begin{compactitem}
							\item This is also going to be a little at a time, over a long period of time.
						\end{compactitem}
				\end{compactitem}
			
		\end{compactitem}
		
		% Why does Moroni want us to remember how merciful God has been throughout time?
		% Quotes DC 19---maybe use the original version?
		% ``And he said unto them: Behold, my bowels are filled with compassion towards you'' (3 Nephi 17:6). Could talk about the ``bowels of mercy''? Connections to OT/NT?
			% Note that 17:7 also talks about it and Christ says it himself.
			% He quotes Alma 7:11--12, which talks about bowels of mercy.
				% Need to talk about the previous verse about the spirit knowing all things (or maybe it's the verse after).
			% Connection to Enoch, Moses 7:41?
			% See my notes at the concepts portion, ``Bowels.''
		% Story of man experiencing healing of his guilt and remorse over causing a serious car accident
		% Key paragraph:
			% ``The Savior’s healing and redeeming power applies to accidental mistakes, poor decisions, challenges, and trials of every kind---as well as to our sins. As I turned to Him, my feelings of guilt and remorse were gradually replaced with peace and rest.''
				% Alma 24:10: ``And I also thank my God, yea, my great God, that he hath granted unto us that we might repent of these things, and also that he hath forgiven us of these our many sins and murders which we have committed and took away the guilt from our hearts through the merits of his Son.''
		% Alma 34:15 > Mosiah 15:8--9 > DC 45:3--5
		% Questions
			% In this talk, Meurs and President Nelson refer to the ``healing and redeeming power'' of Christ. Is there one power to heal, and one power to redeem? Is it all the same power? 
			% What is it about sin that makes it impossible for us to be in God's presence when we have it? 
			% What is the correct way to think about sin: that it causes some cosmic scale to be upset and out of balance; or, that it causes something to in us to be upset and out of balance? Or both?
			% Another way to think about the previous question is, when we receive a remission of our sins, is the remission fixing an out-of-balance cosmic scale, or is it fixing something in us? Or both?
				% When the power of Christ heals us from an ``infirmity,'' or takes away our ``guilt,'' or does anything that isn't fixing a sin, it seems to be fixing something that's in us exclusively, since there is no cosmic scale that could be out of balance. Is it different for sin, then, or not?
		
	% ----------------------
	\subsection{23 April 2023: EQ discussion on the priesthood}
	% ----------------------
		\label{TALKS-23-04-2023}
		
		\begin{compactitem}
			\item Promoting Twitch channel and Substack?
			\item The talk we'll be discussing is Elder Renlund's from the most recent general conference, at the beginning of this month.
				\begin{compactitem}
					\item The talk is called ``Accessing God's power through covenants.''
					\item We'll be reading and talking about certain passages from it, but our goal will be to understand more about some of the things Renlund says about the priesthood.
				\end{compactitem}
			\item Let's begin by reading a passage from Renlund's talk---the sixth paragraph.
				\begin{compactitem}
					\item Here is that paragraph: ``Before the earth was created, God established covenants as the mechanism by which we, His children, could unite ourselves to Him. Based on eternal, unchanging law, He specified the nonnegotiable conditions whereby we are transformed, saved, and exalted. In this life, we make these covenants by participating in priesthood ordinances and promising to do what God asks us to do, and in return, God promises us certain blessings.''
					\item This eternal law is the priesthood, and it gives us the conditions in which we are supposed to perform ordinances and make covenants in order to be ``transformed, saved, and exalted.''
					\item As I was reading about this I was thinking about some of his remarks and some things that the scriptures tell us about the priesthood, or this eternal law.
					\item I'd like to read a bit about that and discuss.
					\item Let's start with \hyperref[DC107-1]{DC 107:1--3}.
					\item In verses 2--3 we learn that the name ``Melchizedek priesthood'' is really just a nickname, and it is not what people called the priesthood before Melchizedek came around.
					\item Before him, they called it ``the holy priesthood, after the order of the Son of God.''
					\item Let's consider a question here: why did they call it ``the holy priesthood, after the order of \textit{the Son of} God''? Why didn't they call it ``the holy priesthood, after the order of God''? Why is it ``the Son of God,'' and not just ``God''?
				\end{compactitem}
			\item As far as I know, the scriptures don't contain a direct answer to this question. We have to look at a number of different things in order to pick up parts of the answer.
				\begin{compactitem}
					\item For example, let's look at \hyperref[JOHN-5-26]{John 5:26} first.
						\begin{compactitem}
							\item Here we read that the Father has ``life in himself,'' and has given, or permitted or allowed or something, the Son to have the same thing, life in himself.
							\item What does it mean to that God has ``life in himself''?
						\end{compactitem}
					\item Now let's read a couple other verses that expand on this idea.
					\item The first will be \hyperref[JOHN-10-17]{John 10:17--18}.
						\begin{compactitem}
							\item Explain and discuss.
						\end{compactitem}
					\item The next will be \hyperref[JOHN-17-2]{John 17:2}.
						\begin{compactitem}
							\item Explain and discuss.
						\end{compactitem}
				\end{compactitem}
			\item Distinction between \gr{ἐξουσία} and \gr{δύναμις}.
				\begin{compactitem}
					\item Explain the distinction
						\begin{compactitem}
							\item \gr{ἐξουσία} means ``power'' in the sense of ``authority'' or ``right to control.''
							\item \gr{δύναμις} means ``power'' in the sense of ``force'' or ``energy,'' a kind of power that \textit{accomplishes} things.
							\item We've seen examples of \gr{ἐξουσία}; there are a few examples of the other kind of power from Christ's life.
							\item Two come from stories you may be familiar with, in \hyperref[LUKE-6-19]{Luke 6:19} and \hyperref[LUKE-8-46]{Luke 8:46}.
							\item Another is \hyperref[LUKE-5-17]{Luke 5:17} and \hyperref[LUKE-5-24]{Luke 5:24}.
						\end{compactitem}
					\item We increase in both as we progress along the covenant path.
					\item Increase in \gr{δύναμις}
						\begin{compactitem}
							\item Perhaps it's like \hyperref[JOHN-4-14]{John 4:14} or \hyperref[DC50-24]{DC 50:24}?
							\item A better example might be \hyperref[1NEPHI-17-48]{1 Nephi 17:48}.
						\end{compactitem}
					\item Increase in \gr{ἐξουσία}
						\begin{compactitem}
							\item We increase in \gr{ἐξουσία} as we make more covenants. 
						\end{compactitem}
					\item Renlund, beginning of the eighth paragraph: ``The term \textit{covenant path} refers to a series of covenants whereby we come to Christ and connect to Him. Through this covenant bond, we have access to His eternal power. The path begins with faith in Jesus Christ and repentance, followed by baptism and receiving the Holy Ghost.''
						\begin{compactitem}
							\item This is an increase in exousia.
							\item See also endnote \#17 in Renlund's talk.
						\end{compactitem}
				\end{compactitem}
			\item Over time, Christ received both more \gr{ἐξουσία} and \gr{δύναμις} as he progressed in his life.
				\begin{compactitem}
					\item We see this in \hyperref[MATTHEW-28-18]{Matthew 28:18}; ``is given'' here is a past tense verb that could be translated ``was given'' or, perhaps even better, ``have been given.''
					\item \hyperref[DC93-12]{DC 93:12--17} appears to be about both types of power.
				\end{compactitem}
			\item So we increase in these ways as we progress---and more specifically, our relationshiph with Jesus Christ changes.
				\begin{compactitem}
					\item Renlund talks about this, in a paragraph about halfway through: ``We become His disciples and represent Him well when we intentionally and incrementally take on ourselves the name of Jesus Christ through covenants. Our covenants give us power to stay on the covenant path because our relationship with Jesus Christ and our Heavenly Father is changed. We are connected to Them by a convenantal bond.''
					\item He says that ``our relationship with Jesus Christ and our Heavenly Father is changed.'' How it is changed? In what way?
					\item We see a description of this after King Benjamin's speech in the book of Mosiah, in chapter 5.
					\item Let's read \hyperref[MOSIAH-5-2]{Mosiah 5:2, 5--6}.
					\item \hyperref[MOSIAH-18-22]{Mosiah 18:22}
					\item \hyperref[MOSIAH-27-25]{Mosiah 27:25}
					\item \hyperref[MOSES-6-64]{Moses 6:64--68}
					\item When we are spiritually born again, we are changed through the power of the atonement---we are given new life through Jesus Christ.
					\item Christ becomes the person who provides us with this new life, new power, and new energy. Christ says, ``Jesus saith unto him, I am the way, the truth, and the life: no man cometh unto the Father, but by me'' (\hyperref[JOHN-14-6]{John 14:6}).
					\item And in this way, he becomes the person through whom we access that power.
					\item Notice that, \textit{for us}, it isn't God that is \textit{our} direct source of this power; it's Christ.
					\item Christ receives it from God, and we receive it from Christ.
					\item And this is why, \textit{for us}, the priesthood is ``after the order of the Son of God,'' because we are integrated into this network of laws and covenants and power in a way where Christ is the source \textit{for us}, even though he isn't the \textit{ultimate} source.
				\end{compactitem}
			\item Return to the question of why the priesthood is after the order of the \textit{Son of God}, rather than just \textit{of God}.
		\end{compactitem}
			
			% XXX % Matthew 28:18: ``And Jesus came and spake unto them, saying, All power is given unto me in heaven and in earth.''
				% Note that ``is given'' is a past tense verb: ``was given'' or ``has been given.'' 
			% XXX % Luke 6:19: ``And the whole multitude sought to touch him: for there went virtue out of him, and healed them all.''
				% δύναμις
			% XXX % Luke 8:46: ``And Jesus said, Somebody hath touched me: for I perceive that virtue is gone out of me.''
				% δύναμιν
			% XXX % John 5:26: ``For as the Father hath life in himself; so hath he given to the Son to have life in himself;''
			% John 6:56--57: ``He that eateth my flesh, and drinketh my blood, dwelleth in me, and I in him. 57 As the living Father hath sent me, and I live by the Father: so he that eateth me, even he shall live by me.''
			% John 10:17--18: ``Therefore doth my Father love me, because I lay down my life, that I might take it again. 18 No man taketh it from me, but I lay it down of myself. I have power to lay it down, and I have power to take it again. This commandment have I received of my Father.
			% XXX % John 14:6: ``Jesus saith unto him, I am the way, the truth, and the life: no man cometh unto the Father, but by me.''
			% XXX % John 17:2: ``As thou hast given him power over all flesh, that he should give eternal life to as many as thou hast given him.''
			% Alma 13:6 (priesthood and atonement)
			% DC 35:2: ``I am Jesus Christ, the Son of God, who was crucified for the sins of the world, even as many as will believe on my name, that they may become the sons of God, even one in me as I am one in the Father, as the Father is one in me, that we may be one.''
			% DC 88:12--13: ``Which light proceedeth forth from the presence of God to fill the immensity of space--13 The light which is in all things, which giveth life to all things, which is the law by which all things are governed, even the power of God who sitteth upon his throne, who is in the bosom of eternity, who is in the midst of all things.''
			% DC 93
				% 2 (important for light)
				% 9 (important for life and light)
				% 16--17 (very important, connects to Matthew 28:18, and talks about the indwelling relationship)
				
	% ----------------------
	\subsection{22 January 2023: As a Child}
	% ----------------------
		\label{TALKS-22-01-2023}
		
		YSA Fifth Ward, Buena Vista, Virgina (note that I didn't actually give this talk on this day; many people were assigned to speak, and there wasn't time. I only shared a brief bit about Mosiah 3:19)
		
		\begin{compactitem}
			\item Introduction to the topic: the new birth in Jesus Christ.
				\begin{compactitem}
					\item I'd like to talk about this a bit by looking at a few verses in scripture.
					\item I'll talk about them and try to make some connections.
				\end{compactitem}
			\item Discuss \hyperref[MOSIAH-3-19]{Mosiah 3:19}, with a focus on whether the list of qualities at the end is meant to explain what it means to become ``as a child.''
				\begin{compactitem}
					\item How much do these qualities actually describe children you know?
					\item Here are some other qualities that children have, which don't appear on this list:
						\begin{compactitem}
							\item Energetic
							\item Curious
							\item Inquisitive
							\item Quick to learn
							\item Quick to remember
							\item Aware that they don't know everything
							\item Willing to let themselves marvel
							\item In love with what is new not because it's new, but because it's an opportunity to learn
							\item Quick to laugh
							\item Playful
							\item Silly
							\item Powerless
							\item Vulnerable
						\end{compactitem}
					\item I suggest that, in our ``becoming as a child,'' gaining or re-gaining these qualities is just as important as gaining the ones we read here that we normally associate with being as a child.
					\item I'd also say that these qualities are far more accurate in describing children than some of the ones we see in the verse.
					\item But we're old now, we're not children anymore. So can we become as a child in this sense?
				\end{compactitem}
			\item Christ is the one that transforms our character so that we can regain, or gain for the first time, these qualities of a child.
				\begin{compactitem}
					\item This is one of the ways to understand the new birth, or the transformatoin that he can carry out in us.
						\begin{compactitem}
							\item \hyperref[MOSIAH-5-7]{Mosiah 5:7}
							\item \hyperref[JOHN-1-12]{John 1:12}
							\item \hyperref[JOHN-3-3]{John 3:3}
							\item \hyperref[ROMANS-12-2]{Romans 12:2}
							\item \hyperref[MOSIAH-27-24]{Mosiah 27:24--26}
							\item \hyperref[ALMA-5-7]{Alma 5:7}
						\end{compactitem}
					\item Christ has the power to carry out this transformation in us.
						\begin{compactitem}
							\item Let's discuss this in \hyperref[JOHN-4-7]{John 4:7--14}.
						\end{compactitem}
					\item This water is a force of energy inside us that will make us \textit{new}. It will make us children again---not in the sense that we become immature or irresponsible or anything like that, but by giving us back those qualities which we might have lost as we've gotten older: energy, curiosity, inquisitiveness, a playful attitude, and so on. We'll get those qualities again and will experience a newness of life that isn't possible any other way.
					\item This newness becomes more valuable the older you get, because the older you get, the farther separated you are from your \textit{actual} childhood.
					\item But that's the power and miracle of the atonement, and that's the power and miracle of Jesus Christ.
					\item The golden age of our lives doesn't have to be in the past, or even in the future. It can be right now, the moments and days we're living \textit{now}, and we can become the kind of person who can make those moments and days a golden age through Christ.
				\end{compactitem}
			\item If I have time, connect to \hyperref[DC84-33]{DC 84:33}.
		\end{compactitem}